\documentclass[11pt]{amsart}
\usepackage{amsmath,amssymb,amsthm,mathtools,mathrsfs}
\usepackage[margin=1.08in]{geometry}
\usepackage{booktabs}
\usepackage{array}
\usepackage{hyperref}
\hypersetup{hidelinks}
\usepackage[final,expansion=false]{microtype}
\setlength{\parskip}{4pt plus 1pt}
\emergencystretch=2.5em

\newtheorem{theorem}{Theorem}[section]
\newtheorem{proposition}[theorem]{Proposition}
\newtheorem{lemma}[theorem]{Lemma}
\newtheorem{corollary}[theorem]{Corollary}
\theoremstyle{definition}
\newtheorem{definition}[theorem]{Definition}
\newtheorem{problem}[theorem]{Problem}
\theoremstyle{remark}
\newtheorem{remark}[theorem]{Remark}

\newcommand{\C}{\mathbb C}
\newcommand{\R}{\mathbb R}
\newcommand{\N}{\mathbb N}
\newcommand{\dd}{\,d}
\newcommand{\RH}{\mathrm{RH}}
\newcommand{\QW}{\mathcal Q_W}
\newcommand{\HL}{\mathcal H_L}
\newcommand{\K}{\mathcal K}
\newcommand{\XiR}{\Xi}
\newcommand{\eps}{\varepsilon}
\DeclareMathOperator{\Spec}{Spec}
\DeclareMathOperator{\Tr}{Tr}
\DeclareMathOperator{\Var}{Var}

\title[Branch selection and global polarization]
{Branch Selection and Global Polarization in the Connes--CCM Route to RH:\\
Finite Prime--Gamma Inertia, Tate Boundary Geometry,\\
and the Divisor-Sensitivity Obstruction}

\author{David Alejandro Trejo Pizzo}
\thanks{Independent Researcher. \texttt{dtrejopizzo@gmail.com}}

\begin{document}

\begin{abstract}
We organize the Connes--Consani--Moscovici and Suzuki operator programs as
a single route from finite semilocal models to the global Weil criterion.
The first part develops a joint prime--Gamma finite-source refinement.  It
keeps the ordinary von Mangoldt channels, the archimedean channel, the
scalar threshold and the exact radical in one Hermitian system.
Haynsworth additivity, scalar Schur innovations, adaptive
Krylov recovery and a heat-core exhaustion reduce the remaining sign, on
every finite row, to one finite inequality (the strict source surplus),
which is not proved; the passage from the finite rows to the Weil form
uses a ground-state identity that is stated as a hypothesis.  For an endpoint-corrected six-mode prolate
source, assuming quoted fixed-order prolate asymptotics, radial
estimates, a lower bound for the interior vector and an extension of the
explicit formula, we bound the Weil form of its exterior co-Poisson
residual by
\(O(\lambda^7d_4)\), and the associated cross functional by
\(O(\lambda^4\sqrt{(1+L)d_4})\), where
\(d_4\asymp\lambda^9e^{-4\pi\lambda^2}\).

The second part constructs the local geometric data that this surplus was
trying to express.  A cyclic root Laplacian has pseudodeterminant \(N^2\);
the half-logarithmic determinant increment
\(\tfrac12\log(\det'\Delta_{mp^k}/\det'\Delta_{mp^{k-1}})\) is \(\log p\),
and its normalized root-stratum incidence is \(p^{-k/2}\).  Their product
is exactly \(\Lambda(p^k)/\sqrt{p^k}\).  At the infinite place, a positive
Gamma spin operator and a two-dimensional polar plane give, by zeta
determinants, the complete factor
\(\frac12s(s-1)\pi^{-s/2}\Gamma(s/2)\).  A shared Tate--Gamma boundary,
matched Abel finite part and covariant charge mixing, with respect to an
explicitly defined source metric, assemble these data
without discarding prime--prime phases.  All these constructions are
source-defined and use neither zeros nor the Weil sign.

We then determine the precise limitation of the ordinary Hilbert
completion.  The complete charged pushout has absolutely continuous scale
spectrum.  The Gamma compliance has dense range by an operator-valued criterion,
and the spectral measures of the source vectors are absolutely
continuous, so the completion has no point spectrum.
Assuming that CCM degree one contains a nonzero eigenclass for each
critical-line zero, every equivariant map from it into this completion
annihilates those classes.  Thus the
local polarization is valid, but this global completion cannot be the
missing polarization of CCM degree one.

Finally, we synthesize six terminal formulations reached by the route---a
signed limit, a Schur surplus, a Rosati factorization, an equivariant
Hilbert majorant, closure-faithful descent, and a positive global
polarization.  For the first, second and last of them the implication to Weil
positivity is proved here (the second only on even tests and under the
ground-state identity, the last under a hypothesis on CCM degree one); the third is a restatement and the fourth and fifth are only
formulated.  In each case the unproved assertion is the terminal
statement itself.  Their common
endpoint is analyzed by two complementary obstructions.  Four forgetful
functors are divisor-blind under explicit deformations (for one of them,
the cyclic Chern character, this is stated, not proved here), while the ordinary
Hilbert pushout fails separately by point-versus-continuous spectral type.
We formulate both mechanisms and isolate two open programs.  One
is geometric: construct an arithmetic intersection product and Hodge-index
polarization on a suitable square of the arithmetic curve.  The other is
analytic but necessarily divisor-sensitive: construct a resonant nuclear
or derived completion retaining evaluation jets before Hilbert reduction.
No proof of either final positivity theorem is claimed.
\end{abstract}

\maketitle
\tableofcontents

\section{Terminal Weil criterion}

Let
\[
 \xi(s)=\frac12s(s-1)\pi^{-s/2}\Gamma(s/2)\zeta(s),
 \qquad
 \XiR(z)=\xi\!\left(\frac12+iz\right).
\]
For a test function \(f\) on \(\R\), let \(W(f)\) denote Weil's explicit
formula distribution, normalized by
\[
\begin{aligned}
 W(f)=&
 \int_\R f(x)(e^{x/2}+e^{-x/2})\,\dd x
 -\sum_{n\ge2}\frac{\Lambda(n)}{\sqrt n}f(\log n)
 -\sum_{n\ge2}\frac{\Lambda(n)}{\sqrt n}f(-\log n)\\
&-(\log 4\pi+\gamma)f(0)
 -\int_0^\infty
 \bigl(f(x)+f(-x)-2e^{-x/2}f(0)\bigr)\frac{e^{x/2}}{e^{x}-e^{-x}}\,\dd x .
\end{aligned}
\tag{1.1}
\]
The associated Hermitian quadratic form is
\[
 \QW(v_1,v_2)=W(v_1*\widetilde v_2),
 \qquad \widetilde v(x)=\overline{v(-x)}.
\tag{1.2}
\]

\begin{theorem}[Weil criterion~\cite{Weil1952,Connes1999}]
\label{thm:weil}
The Riemann Hypothesis is equivalent to
\[
 \QW(v,v)\ge0
 \qquad(v\in C_c^\infty(\R)).
\tag{1.3}
\]
\end{theorem}

\begin{remark}
The route described below does not seek a new final criterion.
The final criterion is (1.3).  Connes' route is a way of constructing
(1.3) as the limit of finite or localized self-adjoint models.
\end{remark}

\begin{remark}[Attribution and novelty boundary]
\label{rem:attribution}
The terminal positivity criterion is Weil's criterion, in the
noncommutative-geometric form developed by Connes.  The Hilbert-space
background around Sonin spaces, co-Poisson summation and radical-type
quotients belongs to Burnol, Connes--Consani,
Connes--Consani--Moscovici, Connes--van Suijlekom and Suzuki
\cite{Burnol2004,Burnol2007,CCArch,CCM2025,CvS2025,Suzuki2025,SuzukiScrew}.
Related Connes--Consani developments around the arithmetic curve and
absolute geometry, and Suzuki's Dirichlet-\(L\) variant, are part of the
same operator-theoretic landscape but are not used as hidden inputs below
\cite{CCJacobian,CCAbsolute,SuzukiDirichlet}.
The contribution of this paper starts where these inputs leave an unsigned branch-selection
problem: finite-source Schur inertia, prime-by-prime innovation,
observability staircases, adaptive Krylov recovery, the six-mode prolate
residual estimates, and the heat/hybrid form-core exhaustion.  Classical
linear-algebraic devices used below, such as shorted operators,
Haynsworth inertia additivity, Loewner order and Andreief's identity, are
cited at their first structural use.
\end{remark}

\section{The CCM finite spectral chain}

Put \(L=2\log\lambda\).  The semilocal construction restricts the Weil form
to a finite Fourier space
\[
 V_{L,N}=\operatorname{span}\{e^{2\pi ikx/L}: |k|\le N\}.
\tag{2.1}
\]
Let \(W_{L,N}\) be the corresponding finite Weil matrix.  It contains the
polar, archimedean and prime-power terms of the semilocal trace formula.
The finite structural facts used below are those of the CCM/Connes and Connes--van Suijlekom finite programs
\cite{CCM2025,Connes2026,CvS2025}.

\begin{definition}[Finite CCM package]
\label{def:ccm-package}
For each \(L,N\), the finite package consists of:
\[
 \eps_{L,N}=\min\Spec W_{L,N},
 \qquad
 T_{L,N}=W_{L,N}-\eps_{L,N}I,
\tag{2.2}
\]
the radical line \(\C\xi_{L,N}=\ker T_{L,N}\) when the bottom is simple,
and the positive quotient
\[
 \K_{L,N}=(\C\xi_{L,N})^\perp,
 \qquad
 S_{L,N}=T_{L,N}|_{\K_{L,N}}.
\tag{2.3}
\]
\end{definition}

Connes--Consani--Moscovici identify the finite commutator structure,
construct the quotient operator, and obtain self-adjoint finite spectral
data after shifting and quotienting~\cite{CCM2025}.  Connes' 2026
extremal construction
adds the prolate model \(k_L\), whose transform is proposed there as an approximation of \(\XiR\) (see Remark~\ref{rem:prolate-limit}).
The finite spectral statements used below are the following.

\begin{theorem}[Finite critical-line CCM--Connes model~\cite{CCM2025,Connes2026,CvS2025}]
\label{thm:ccm-input}
For each finite level \((L,N)\) at which the lowest eigenvalue of
\(W_{L,N}\) is simple with even eigenvector, the quotient operator
\(S_{L,N}\) is self-adjoint in the positive quotient metric, and the
zeros of the Fourier transform of the corresponding finite ground vector
are real, that is, its finite spectral divisor lies on the critical line.
\end{theorem}

\begin{remark}
\label{rem:prolate-limit}
Theorem~\ref{thm:ccm-input} is quoted from the cited works and is not
reproved here; we do not attach statement numbers to it, and the reader
should consult \cite{CCM2025,CvS2025} for the precise hypotheses under
which the reality of the finite zeros is established.  In addition,
Connes~\cite{Connes2026} describes an explicit prolate model \(k_L\)
and discusses the approximation
\[
 \widehat k_L(z)\longrightarrow \XiR(z)
\tag{2.4}
\]
locally uniformly on compact subsets of the open strip
\(|\operatorname{Im}z|<1/2\).  We do not use (2.4) as an established
input: every statement below that involves it says ``assume (2.4)''.
What these finite statements do not give is lower-semicontinuity of the
indefinite Weil form in the global limit.
\end{remark}

\section{Localized operator realization}

Suzuki replaces the distributional Weil form by a continuous screw-function
coordinate.  Let \(g\) be the screw function associated with zeta.  On
\((-a,a)\) one obtains a localized quadratic form \({\QW}_a\) and a
self-adjoint operator \(A_a\) associated with it.  The nonlocal
first-order realization provides a concrete operator-theoretic
interpretation of the finite intervals.

\begin{theorem}[Localized self-adjoint Suzuki realization~\cite{Suzuki2025,SuzukiScrew}]
\label{thm:suzuki-input}
For each finite interval \((-a,a)\), the localized Weil quadratic form has
a self-adjoint operator realization \(A_a\), as constructed in
\cite{Suzuki2025,SuzukiScrew}; see there for the range of \(a\) and the
hypotheses.
\end{theorem}

\begin{remark}
Suzuki's limiting conjecture states that the imaginary parts of the
nontrivial zeros arise from the strong-resolvent limit of these
finite-interval nonlocal realizations as \(a\to\infty\).
Theorem~\ref{thm:suzuki-input} supplies the correct analytic language for convergence: strong
resolvent convergence of finite-interval operators.  It does not prove the
limit that would put the full zero divisor on the line.
\end{remark}

\section{Finite-to-global branch convergence}

There are two different statements which must not be identified.

\begin{definition}[Divisor convergence]
Let \(\phi_L\) be the selected continuum ground branch obtained from
the finite radical quotients after \(N\to\infty\) at fixed \(L\),
whenever this limit exists.  Divisor convergence means that each
\(\widehat\phi_L\) is holomorphic and has only real zeros in the strip
\(|\operatorname{Im}z|<1/2\), and that, locally uniformly on compact
subsets of that strip avoiding the zeros of \(\XiR\),
\[
 \partial_z^2\log \widehat\phi_L(z)
 \longrightarrow
 \partial_z^2\log \XiR(z).
\tag{4.1}
\]
\end{definition}

\begin{definition}[Positive-form convergence]
Positive-form convergence means that for every \(v\in C_c^\infty(\R)\)
there are finite-level approximants \(v_{L,N}\in V_{L,N}\) with
\[
 v_{L,N}\to v
\tag{4.2}
\]
in a topology strong enough for the explicit formula, and
\[
 \QW(v,v)
 \ge
 \limsup_{L\to\infty}\limsup_{N\to\infty}
 \Bigl\langle T_{L,N}v_{L,N},v_{L,N}\Bigr\rangle .
\tag{4.3}
\]
\end{definition}

The finite forms on the right side are nonnegative after shift and
quotient.  Thus positive-form convergence immediately implies Weil
positivity.  Divisor convergence also implies RH, by the argument in
the proof of Theorem~\ref{thm:closure}.

\begin{lemma}[Logarithmic-curvature branch form]
\label{lem:branch}
Assume (2.4).  Let \(k_L\) be Connes' prolate model and \(\phi_L\) the
selected ground branch.  Then the convergence (4.1) is equivalent, on
every simply connected compact subset of the strip
\(|\operatorname{Im}z|<1/2\) avoiding the zeros of \(\XiR\) and of the
\(\widehat k_L\), to
\[
 \boxed{
 \partial_z^2\log
 \frac{\widehat\phi_L(z)}{\widehat k_L(z)}
 \longrightarrow0 .}
\tag{4.4}
\]
\end{lemma}

\begin{proof}
Choose compatible logarithms on the compact set.  Then
\[
 \partial_z^2\log\widehat\phi_L-\partial_z^2\log\XiR
 =
 \partial_z^2\log\frac{\widehat\phi_L}{\widehat k_L}
 +
 \partial_z^2\log\frac{\widehat k_L}{\XiR}.
\tag{4.5}
\]
The last term tends to zero by local uniform convergence and Cauchy
estimates.  This proves the equivalence.
\end{proof}

\section{Conditional closure of the CCM route}

\begin{theorem}[Conditional closure of the CCM route]
\label{thm:closure}
Assume the finite CCM package.  If either positive-form convergence or
divisor convergence holds, then RH follows.
\end{theorem}

\begin{proof}
Under positive-form convergence, the finite shifted quotient forms are
nonnegative:
\[
 \langle T_{L,N}u,u\rangle\ge0
 \qquad(u\perp\xi_{L,N}).
\tag{5.1}
\]
Given \(v\in C_c^\infty(\R)\), choose approximants \(v_{L,N}\) as in
(4.2)--(4.3).  Passing to the limit gives
\[
 \QW(v,v)\ge0.
\tag{5.2}
\]
By Theorem~\ref{thm:weil}, RH follows.

Under divisor convergence, recall that every zero of \(\XiR\) lies in
the strip \(|\operatorname{Im}z|<1/2\), since \(\zeta\) has no zeros on
\(\operatorname{Re}s=0,1\).  Suppose \(z_0\) is a zero of \(\XiR\) of
order \(m\ge1\) with \(\operatorname{Im}z_0\ne0\).  Choose
\(0<r<|\operatorname{Im}z_0|\) so that the closed disc
\(\overline D=\{|z-z_0|\le r\}\) lies in the strip and contains no other
zero of \(\XiR\).  Since the zeros of \(\widehat\phi_L\) in the strip
are real, \(F_L=\partial_z^2\log\widehat\phi_L\) is holomorphic on a
neighbourhood of \(\overline D\).  By (4.1), \(F_L\) converges uniformly
on the annulus \(A=\{r/2\le|z-z_0|\le r\}\).  By the maximum principle,
\((F_L)\) is uniformly Cauchy on \(\overline D\), hence converges
uniformly there to a holomorphic function \(F\).  On \(A\),
\(F=\partial_z^2\log\XiR\), so by the identity theorem
\(\partial_z^2\log\XiR\) extends holomorphically to \(z_0\).  But
\(\partial_z^2\log\XiR(z)=-m(z-z_0)^{-2}+O(1)\) near \(z_0\), a
contradiction.  Hence every zero of \(\XiR\) is real, that is, all
zeros of \(\xi\) lie on the critical line.  This is RH.
\end{proof}

\begin{remark}
Theorem~\ref{thm:closure} is conditional: neither form of convergence is
proved here.  Everything finite is self-adjoint and critical-line.  The only possible failure is at
the infinite passage: positivity or spectral reality may be lost if the
finite branch does not converge to the correct global branch with sign.
\end{remark}

\section{Localized realization and the global limit}

Suzuki's screw-function operator provides the natural candidate topology for
the limiting passage.  In the notation of Theorem~\ref{thm:suzuki-input},
the desired operator statement is
\[
 A_a \longrightarrow A_\infty
\tag{6.1}
\]
in strong resolvent sense, with
\[
 \Spec(A_\infty)=\{\gamma:\xi(1/2+i\gamma)=0\}
\tag{6.2}
\]
including multiplicities.

\begin{proposition}[Suzuki-to-Connes transfer]
Suppose that (i) there are isometries \(J_L\) of the CCM quotient spaces
into one Hilbert space \(\mathcal H\) such that
\(J_LS_LJ_L^*\to A_\infty\) in the strong resolvent sense, where
\(S_L\) denotes the quotient operators along a cofinal diagonal and
\(A_\infty\) is a densely defined operator as in (6.1); and (ii)
\(\Spec(A_\infty)\), with multiplicities, is
\(\{\gamma:\xi(1/2+i\gamma)=0\}\) as in (6.2).  Then all zeros of
\(\xi\) lie on the critical line.
\end{proposition}

\begin{proof}
A strong resolvent limit of self-adjoint operators which is a densely
defined operator is self-adjoint, so its spectrum is real.  By (ii) every
\(\gamma\) with \(\xi(1/2+i\gamma)=0\) is real.
\end{proof}

\begin{remark}
Hypothesis (ii), for any self-adjoint \(A_\infty\) (for instance the
Suzuki limit (6.1), if it exists), already gives RH; hypothesis (i) is
used only to make \(A_\infty\) self-adjoint.  The proposition therefore
only records where the identification of the limiting spectrum with the
zero divisor enters; it does not yield the divisor convergence (4.1) for
the CCM branch \(\phi_L\).  The branch-selection step can thus be posed
as branch convergence of transforms, as form convergence, or as
strong-resolvent convergence together with spectral identification.
\end{remark}

\section{Branch selection as a finite source theorem}

The finite critical-line property and the prolate approximation do not
alone prove RH.  The branch-selection assertion is a signed convergence
theorem.  In
the notation above it can be stated in three sufficient operational forms.

\begin{problem}[Branch selector]
\label{prob:branch}
Assume (2.4).  Let \(\phi_L\) be the ground branch selected by the
finite CCM radical quotient after \(N\to\infty\), and assume that it
exists and that \(\widehat\phi_L\) is holomorphic with only real zeros
in \(|\operatorname{Im}z|<1/2\).  Prove that
\[
 \partial_z^2\log
 \frac{\widehat\phi_L(z)}{\widehat k_L(z)}
 \longrightarrow0
\tag{7.1}
\]
locally uniformly on compact subsets of the open critical strip avoiding
zeros.
\end{problem}

\begin{problem}[Positive finite-to-global transfer]
\label{prob:transfer}
For every \(v\in C_c^\infty(\R)\), construct finite quotient approximants
\(v_{L,N}\) for which the completed Weil form satisfies
\[
 \QW(v,v)
 =
 \lim_{L\to\infty}\lim_{N\to\infty}
 \langle T_{L,N}v_{L,N},v_{L,N}\rangle
 +
 R(v),
\tag{7.2}
\]
with \(R(v)\ge0\), or at least with \(R(v)\) controlled so that the right
side is nonnegative.
\end{problem}

\begin{problem}[All-source sharp Doob floor]
\label{prob:doob}
Let \(K\) be Riemann's full kernel, normalized by
\(\widehat K=\XiR\); it is even and strictly positive
(P\'olya~\cite{Polya1926}).  Set
\[
 d\mu_K(x)=2\cosh(x/2)K(x)\,\dd x .
\tag{7.3}
\]
For even compactly supported smooth \(r\) define the jump energy
\[
 \mathscr E_K(r)=\QW(Kr,Kr)+\frac12\Var_{\mu_K}(r),
\]
so that \(\mathscr E_K(r)-\frac12\Var_{\mu_K}(r)=\QW(Kr,Kr)\) holds by
definition.  Prove, directly from the explicit formula (1.1), that is,
from the von Mangoldt atoms, the archimedean term and the pole, the sharp
inequality
\[
 \boxed{
 \mathscr E_K(r)\ge\frac12\Var_{\mu_K}(r)}
\tag{7.4}
\]
for every even compactly supported smooth \(r\).
\end{problem}

\begin{proposition}[What each gate gives]
Each of Problems~\ref{prob:branch} (under its stated assumptions) and
\ref{prob:transfer} implies RH.  Problem~\ref{prob:doob} is equivalent to
Weil positivity on even tests; together with the corresponding inequality
for odd tests it gives (1.3), hence RH.
\end{proposition}

\begin{proof}
Problem~\ref{prob:branch} gives (4.1) by Lemma~\ref{lem:branch}; with
the assumed existence and real-zero property of \(\phi_L\) this is
divisor convergence, hence RH by Theorem~\ref{thm:closure}.
Problem~\ref{prob:transfer} gives Weil positivity directly, hence RH by
Theorem~\ref{thm:weil}.  For Problem~\ref{prob:doob}, the definition of
\(\mathscr E_K\) gives
\[
 \QW(Kr,Kr)=\mathscr E_K(r)-\frac12\Var_{\mu_K}(r).
\tag{7.5}
\]
Since \(K>0\), every even compactly supported smooth test has the form
\(Kr\) with \(r\) even, compactly supported and smooth.  Hence (7.4)
is the Weil criterion (1.3) restricted to even tests.  Since \(W\) is
even, \(\QW(v,v)=\QW(v_e,v_e)+\QW(v_o,v_o)\) for the even and odd parts
\(v_e,v_o\) of \(v\), so positivity on odd tests in addition gives (1.3)
for all tests, hence RH by Theorem~\ref{thm:weil}.  Whether even
positivity alone implies RH is not decided here.  (Under RH, Weil
positivity holds and (7.4) follows from (7.5); the converse for the other
two gates is not claimed here.)
\end{proof}

\begin{remark}
With the definition above, Problem~\ref{prob:doob} is a restatement of
even Weil positivity.  Its intended content is that \(\mathscr E_K(r)\)
be expanded, by a ground-state transform based on
\(\XiR(\pm i/2)=\tfrac12\) and \(\int d\mu_K=1\), as a manifest jump form over
the von Mangoldt atoms and the archimedean jump density; such an
expansion is not derived in this paper.
\end{remark}

\section{Finite branch-selection refinement}

The gate which belongs to the Connes route is not the full-kernel
inequality (7.4), because that has already eliminated the finite
semilocal architecture.  The required finite comparison is instead a
one-sided Mosco
transfer for the finite shifted forms, in the sense of Mosco convergence
of closed forms~\cite{Mosco1969,KuwaeShioya2003,Post2012}.  We use Mosco
only as a directional lower-transfer requirement; known instability
phenomena for global path properties under Mosco convergence are a
warning not to infer additional global structure from this step alone
\cite{SuzukiUemura2014}.

Fix \(v\in C_c^\infty(\R)\).  Let \(P_{L,N}v\in V_{L,N}\) be the Fourier
sampling used in the CCM finite model, and remove the finite radical by
\[
 u_{L,N}(v)
 =
 P_{L,N}v-
 \frac{\langle P_{L,N}v,\xi_{L,N}\rangle}
      {\langle \xi_{L,N},\xi_{L,N}\rangle}\xi_{L,N}.
\tag{8.1}
\]
Every object in (8.1) is finite and is defined before any reference to
the zero locations of \(\zeta\).

\begin{definition}[Finite transfer defect]
The finite transfer defect of \(v\) is
\[
 \mathfrak D_{L,N}(v)
 =
 \QW(v,v)-
 \langle T_{L,N}u_{L,N}(v),u_{L,N}(v)\rangle .
\tag{8.2}
\]
Expanding both terms by the explicit formula and by the semilocal trace
formula decomposes it as
\[
 \mathfrak D_{L,N}
 =
 \mathfrak D^{\rm mesh}_{L,N}
 +\mathfrak D^{\rm ext}_{L,N}
 +\mathfrak D^{\Gamma}_{L,N}
 +\mathfrak D^{p}_{L,N}
 +\mathfrak D^{\rm pole}_{L,N}
 +\mathfrak D^{\rm br}_{L,N}.
\tag{8.3}
\]
Here \(\mathfrak D^{\rm mesh}\) is Fourier discretization, 
\(\mathfrak D^{\rm ext}\) is the exterior lattice factor,
\(\mathfrak D^{\Gamma}\), \(\mathfrak D^p\), and
\(\mathfrak D^{\rm pole}\) are the literal Gamma, von Mangoldt and polar
source differences, and \(\mathfrak D^{\rm br}\) is the finite radical
branch correction.  These six terms are not given by separate formulas
here: (8.3) is a bookkeeping convention, and the terms are fixed only up
to the constraint that they sum to (8.2).
\end{definition}

\begin{proposition}[Analytic part of the transfer, conditional]
\label{prop:known-transfer}
Fix \(v\in C_c^\infty(\R)\) and assume that the decomposition (8.3)
has been fixed so that
\(|\mathfrak D^{\rm ext}_{L,N}(v)|\le C_vL^2/N\) and
\(\mathfrak D^{\rm mesh}_{L,N}(v)\to0\) as \(N\to\infty\) at fixed
\(L\).  Then there is a cofinal diagonal \(N=N(L)\), with
\(N(L)/L^2\to\infty\), such that
\[
 \mathfrak D^{\rm mesh}_{L,N(L)}(v)
 +\mathfrak D^{\rm ext}_{L,N(L)}(v)
 \longrightarrow0 .
\tag{8.4}
\]
\end{proposition}

\begin{proof}
For each \(L\) choose \(N(L)\ge L^3+L^2\) so large that
\(|\mathfrak D^{\rm mesh}_{L,N(L)}(v)|\le1/L\).  Then
\(N(L)/L^2\to\infty\) and
\(|\mathfrak D^{\rm ext}_{L,N(L)}(v)|\le C_vL^2/N(L)\to0\).
\end{proof}

\begin{problem}[Source-side Mosco lower bound]
\label{prob:source-mosco}
Prove that for every \(v\in C_c^\infty(\R)\) one can choose the cofinal
diagonal of Proposition~\ref{prop:known-transfer} so that
\[
 \liminf_{L\to\infty}
 \Bigl(
 \mathfrak D^{\Gamma}_{L,N(L)}(v)
 +\mathfrak D^{p}_{L,N(L)}(v)
 +\mathfrak D^{\rm pole}_{L,N(L)}(v)
 +\mathfrak D^{\rm br}_{L,N(L)}(v)
 \Bigr)
 \ge0 .
\tag{8.5}
\]
\end{problem}

\begin{theorem}[Mosco lower transfer implies RH]
\label{thm:mosco-transfer}
Assume the hypotheses of Proposition~\ref{prop:known-transfer} for every
\(v\in C_c^\infty(\R)\).  If Problem~\ref{prob:source-mosco} holds,
then the positive finite CCM
forms transfer to the completed Weil form.  Consequently RH follows.
\end{theorem}

\begin{proof}
For the diagonal in Proposition~\ref{prop:known-transfer} and
Problem~\ref{prob:source-mosco}, equations (8.2)--(8.5) give
\[
 \mathfrak D_{L,N(L)}(v)\ge-o(1).
\tag{8.6}
\]
Since \(T_{L,N}\ge0\) on the finite quotient, (8.2) implies
\[
 \QW(v,v)
 =
 \langle T_{L,N(L)}u_{L,N(L)}(v),u_{L,N(L)}(v)\rangle
 +\mathfrak D_{L,N(L)}(v)
 \ge-o(1).
\tag{8.7}
\]
Letting \(L\to\infty\) gives \(\QW(v,v)\ge0\).  Since \(v\) was
arbitrary, Weil positivity follows.  Theorem~\ref{thm:weil} then gives
RH.
\end{proof}

\begin{remark}
Problem~\ref{prob:source-mosco} is not the full-kernel inequality (7.4)
but a finite-source statement: the Gamma term, the von Mangoldt atoms,
the pole and the radical correction are kept together until after the
lower limit.  Bounding them separately does not give a sign.
\end{remark}

\begin{remark}[Joint-source rule]
The source-side lower bound is studied below through the complete finite
prime--Gamma form and its radical anti-short.  The archimedean, polar,
ordinary-prime and radical contributions may exchange mass, so no
separate lower bound is imposed on any one of them.  The finite
statement used is the joint inertia/gain inequality of
Section~\ref{sec:cofinal-inertia}.
\end{remark}

\subsection{Controlled exterior and cross residuals}

We now specify the moving vector and bound its even Rayleigh residual.
This specification is essential: a trial
vector satisfying only the two co-Poisson moment conditions may retain a
nonzero endpoint value and hence a literal \(t^{-1}\) exterior carrier.
The following fixed six-mode construction also removes the first two
exterior boundary symbols without changing the \(d_4,d_8\) angle ladder.

Let
\[
 J_6=\{0,4,8,12,16,20\}
\]
and let \(\psi_{j,\lambda}\) be the real normalized even prolate
spheroidal wave functions in the Fourier \(+1\) sector.  Write
\[
 a_j=\psi_{j,\lambda}(0),\qquad
 e_j=\psi_{j,\lambda}(\lambda),\qquad
 q_j=1-\chi_j,\qquad d_j=1-\chi_j^2,
\tag{8.8a}
\]
where \(\chi_j\) is the corresponding finite Fourier eigenvalue, and let
\(\mu_j\) be the separation constant in the prolate differential equation.
Here \(\mathcal F f(t)=\int_\R f(x)e^{-2\pi ixt}\,\dd x\);
\(\psi_{j,\lambda}\) is the \(j\)-th eigenfunction of
\(P_\lambda\mathcal F P_\lambda\) on \(L^2(-\lambda,\lambda)\), of unit
norm there and extended by zero, with
\(P_\lambda\mathcal F\psi_{j,\lambda}=\chi_j\psi_{j,\lambda}\) and
\(0<\chi_j<1\) for \(j\equiv0\bmod4\) (\(P_\lambda\) is defined below);
it satisfies
\(((\lambda^2-t^2)\psi')'-(2\pi\lambda t)^2\psi+\mu_j\psi=0\) on
\((-\lambda,\lambda)\).  Throughout, \(\widehat f=\mathcal F f\).

\emph{Co-Poisson vectors.}  Let \(P_\lambda\) be multiplication by
\(1_{[-\lambda,\lambda]}\).  For \(f=\sum_{j\in J_6}c_j\psi_{j,\lambda}\) put
\(G_f(v)=v^{1/2}\sum_{n\ge1}\widehat f(nv)\) for \(v>0\), and, for
\(v\ge\lambda\), \(H_f(v)=v^{1/2}\sum_{n\ge1}((1-P_\lambda)\widehat f)(nv)\),
which equals \(G_f(v)\) there.  In the variable \(x=\log v\), with
\(I_L=[-L/2,L/2]\) and \(\lambda=e^{L/2}\), set
\[
\begin{gathered}
 V_f^+(x)=\frac{g(x)+g(-x)}{\sqrt2},\quad g=G_f(e^x)1_{I_L}(x);\\
 R_f^+(x)=\frac{h(x)+h(-x)}{\sqrt2},\quad h=H_f(e^x)1_{[L/2,\infty)}(x).
\end{gathered}
\]
Thus \(V_f^+\in L^2(I_L)_{\rm even}\) is the interior co-Poisson vector
and \(R_f^+\), supported in \(\R\setminus(-L/2,L/2)\), is the exterior
residual.  For even functions \(u\) we write
\(\widehat u(w)=\int_\R u(x)e^{-wx}\,dx\) for the Mellin coordinate
wherever the integral converges.

\emph{Standing assumptions.}  The estimates of this subsection are
derived under the following inputs, which are quoted from the prolate
literature or assumed, and are not proved here.
\begin{itemize}
\item[(P1)] The fixed-order prolate asymptotics (8.8h)--(8.8i) and
(8.8au) (Fuchs~\cite{Fuchs1964}, Slepian~\cite{Slepian1965}: for fixed
\(n\) and \(c=2\pi\lambda^2\to\infty\), one has
\(1-\chi_n^2\sim\pi^{1/2}2^{3n+2}c^{n+1/2}e^{-2c}/n!\)), together with
\(e_j^2\asymp\lambda d_j\) for \(j\in J_6\).
\item[(P2)] The radial estimates (8.8n), (8.8q)--(8.8r) (with
\(U_{0,\lambda},U_{1,\lambda}\) independent of \(j\)) and (8.8v), for the
ranges \(1\le x\le2\), \(2\le x\le c\) and \(x\ge c\) respectively,
as derived from the expansions of Dunster~\cite{DunsterPSWF}.
\item[(P3)] The lower bound (8.8e) for the interior vector.
\item[(P4)] The explicit formula (1.1)--(1.2) extends, with absolute
convergence and symmetric summation over zeros, to even, piecewise
\(C^1\) functions \(u\) with finitely many jumps and
\(|u(x)|+|u'(x)|\ll e^{-(1/2+\delta)|x|}\) for some \(\delta>0\), in the form \(\QW(u,u)=\sum_\rho\widehat u(\rho-\tfrac12)
\overline{\widehat u(\bar\rho-\tfrac12)}\), and to the sesquilinear
pairing of such a \(u\) with \(g\in L^2(I_L)\).
\end{itemize}
For the vector \(f_\lambda^{\rm ep}\) of
Definition~\ref{def:endpoint-corrected-source} we write
\(V_\lambda^+,R_\lambda^+\) for \(V_f^+,R_f^+\), and we define the
even Rayleigh residual and the cross functional by
\[
 \mathscr R_L^+=\frac{\QW(R_\lambda^+,R_\lambda^+)}{\|V_\lambda^+\|^2},
 \qquad
 \|b_L^+\|=\sup\Bigl\{\Bigl|\QW\Bigl(\frac{R_\lambda^+}{\|V_\lambda^+\|},g\Bigr)\Bigr|:
 g\in L^2(I_L)_{\rm even},\ \|g\|=1,\ g\perp V_\lambda^+\Bigr\}.
\]

\begin{definition}[Endpoint-corrected moving source]
\label{def:endpoint-corrected-source}
Set
\[
 \mathcal K_\lambda^{\rm ep}
 =\left\{
 f=\sum_{j\in J_6}c_j\psi_{j,\lambda}:
 f(0)=0,\quad \int_{\mathbb R}f=0,\quad
 f(\lambda)=f'(\lambda)=0
 \right\}.
\tag{8.8b}
\]
Here \(f(\lambda),f'(\lambda)\) are the limits from inside
\((-\lambda,\lambda)\).
Let \(f_\lambda^{\rm ep}\) be a normalized first generalized eigenvector
of \(\sum d_j|c_j|^2\) relative to \(\sum|c_j|^2\) on this space, and let
\(V_\lambda^+=V^+_{f_\lambda^{\rm ep}}\) and
\(R_\lambda^+=R^+_{f_\lambda^{\rm ep}}\) be its interior and exterior
co-Poisson vectors.  Henceforth the moving vector is
\[
 q_L^+=\frac{V_\lambda^+}{\|V_\lambda^+\|}.
\tag{8.8c}
\]
This is also the vector used in the cross residual of
Theorem~\ref{thm:cross-residual-vanishing}.
\end{definition}

\begin{lemma}[The endpoint constraints preserve the angle ladder]
\label{lem:endpoint-angle-ladder}
Assume (P1).  For all sufficiently large \(\lambda\), the four
constraints in
(8.8b) are independent.  If
\(\delta_{0,\lambda}^{\rm ep}\le
\delta_{1,\lambda}^{\rm ep}\) are the generalized eigenvalues of
\(\sum d_j|c_j|^2\) relative to \(\sum|c_j|^2\) on
\(\mathcal K_\lambda^{\rm ep}\), then
\[
 \delta_{0,\lambda}^{\rm ep}\asymp d_4,
 \qquad
 \delta_{1,\lambda}^{\rm ep}\asymp d_8.
\tag{8.8d}
\]
The lower bound
\[
 \|V_\lambda^+\|\ge c_V>0,
\tag{8.8e}
\]
with \(c_V\) independent of \(\lambda\), is not proved here; it is the
standing assumption (P3).
\end{lemma}

\begin{proof}
The exact endpoint-slope identity for a prolate linear combination is
\[
 2\lambda f'(\lambda)
 =\sum_{j\in J_6}c_j(\mu_j-4\pi^2\lambda^4)e_j.
\tag{8.8f}
\]
After imposing \(f(\lambda)=0\), the four constraint rows are therefore
\[
 (a_j),\qquad(q_ja_j),\qquad(e_j),\qquad(\mu_je_j).
\tag{8.8g}
\]
For \(j=4r\), \(0\le r\le5\), the fixed-order prolate-to-Hermite and
endpoint asymptotics of (P1) give
\[
 |a_{4r}|\asymp1,\qquad
 |e_{4r}|\asymp\lambda^{1/2}q_{4r}^{1/2},\qquad
 \mu_{4r}=\lambda^2(M_0+M_1r+O(\lambda^{-2})),
\tag{8.8h}
\]
where \(M_1\ne0\), together with
\[
 q_{4(r+1)}/q_{4r}\asymp\lambda^8.
\tag{8.8i}
\]
Put \(y=\lambda^4\) and use the unknowns \(w_r=c_{4r}e_{4r}\).  Since
\(\chi_j\to1\), \(q_j\le d_j\le2q_j\), and (8.8h)--(8.8i) give
\(q_{4r}\asymp q_0y^{2r}\).  Subtracting \(\lambda^2M_0\) times the third
row from the fourth and dividing by \(\lambda^2M_1\), the constraints
become \(Mw=0\) with
\[
 M=\begin{pmatrix}
 \beta_0&\beta_1y^{-1}&\beta_2y^{-2}&\cdots&\beta_5y^{-5}\\
 \gamma_0&\gamma_1y&\gamma_2y^2&\cdots&\gamma_5y^5\\
 1&1&1&\cdots&1\\
 \epsilon_0&1+\epsilon_1&2+\epsilon_2&\cdots&5+\epsilon_5
 \end{pmatrix},
\tag{8.8j}
\]
where \(|\beta_r|\asymp|\gamma_r|\asymp1\) uniformly (after dividing the
first two rows by \(|a_0/e_0|\) and \(|q_0a_0/e_0|\)) and
\(\epsilon_r=O(\lambda^{-2})\).  Write \(\|w\|=\max_r|w_r|\).  Every
\(w\in\ker M\) satisfies, from the first two rows,
\[
 |w_0|\ll y^{-1}\|w\|,\qquad |w_5|\ll y^{-1}\|w\|.
\tag{8.8k}
\]

\emph{Independence.}  If \(\dim\ker M\ge3\), it contains \(w\ne0\) with
\(w_1=w_2=0\).  The last two rows then give
\(w_3+w_4=O(y^{-1}\|w\|)\) and
\(3w_3+4w_4=O((y^{-1}+\lambda^{-2})\|w\|)\), so
\(\|w\|\ll(y^{-1}+\lambda^{-2})\|w\|\), impossible for large \(\lambda\).
Hence \(M\) has rank four and \(\ker M\) is two-dimensional.

\emph{The quotient.}  Since \(e_{4r}^2\asymp\lambda d_{4r}\) and
\(d_{4r}\asymp d_4y^{2(r-1)}\), for \(w\in\ker M\)
\[
 \frac{\sum_jd_j|c_j|^2}{\sum_j|c_j|^2}
 \asymp\frac{\sum_r|w_r|^2}{\sum_r|w_r|^2/d_{4r}} .
\]
By (8.8k), \(|w_0|^2/d_0\ll\|w\|^2/d_4\), and
\(|w_r|^2/d_{4r}\le\|w\|^2/d_4\) for \(r\ge1\); hence every kernel vector
has quotient \(\gg d_4\), so \(\delta_{0,\lambda}^{\rm ep}\gg d_4\).
Conversely, take \(w\in\ker M\setminus\{0\}\) with \(w_5=0\).  If
\(|w_1|\le\eta\|w\|\) with \(\eta\) small, the second row gives
\(|w_4|\ll y^{-1}\|w\|\), and then the last two rows force
\(|w_2|+|w_3|\ll(\eta+y^{-1}+\lambda^{-2})\|w\|\), a contradiction.  So
\(|w_1|\gg\|w\|\), the denominator is \(\gg\|w\|^2/d_4\), and the
quotient is \(\ll d_4\).  Thus \(\delta_{0,\lambda}^{\rm ep}\asymp d_4\).

Since \(\ker M\) is two-dimensional, \(\delta_{1,\lambda}^{\rm ep}\) is
the maximum of the quotient on \(\ker M\).  For any kernel vector, if
\(\max_{r\le2}|w_r|\le\eta\|w\|\), then \(|w_5|\ll y^{-1}\|w\|\) and the
last two rows force \(|w_3|+|w_4|\ll(\eta+y^{-1}+\lambda^{-2})\|w\|\), a
contradiction; hence \(\max_{r\le2}|w_r|\gg\|w\|\), the denominator is
\(\gg\|w\|^2/d_8\), and the quotient is \(\ll d_8\).  Conversely, take
\(w\in\ker M\setminus\{0\}\) with \(w_0=0\).  The first row gives
\(|w_1|\ll y^{-1}\|w\|\), so \(|w_1|^2/d_4\ll\|w\|^2/d_8\), while
\(|w_r|^2/d_{4r}\le\|w\|^2/d_8\) for \(r\ge2\); the quotient is
\(\gg d_8\).  This proves (8.8d).
\end{proof}

Put
\[
 \ell_\lambda=(1-P_\lambda)\widehat f_\lambda^{\rm ep}.
\tag{8.8l}
\]

\begin{lemma}[Uniform exterior two-jet estimate]
\label{lem:uniform-exterior-two-jet}
Assume (P1) and (P2).  Uniformly for \(t\ge\lambda\),
\[
 \boxed{
 |\ell_\lambda(t)|+
 \lambda^{-1}|\ell_\lambda'(t)|
 \ll \lambda^{7/2}\sqrt{d_4}\,t^{-3}.}
\tag{8.8m}
\]
\end{lemma}

\begin{proof}
Put \(c=2\pi\lambda^2\) and \(x=t/\lambda\).  The Fourier continuation
of each fixed prolate mode is its radial PSWF.  The estimates (8.8n),
(8.8q)--(8.8r) and (8.8v) below are those of assumption (P2); they are
what the uniform Bessel and Liouville--Green expansions of
Dunster~\cite{DunsterPSWF}, with their differentiated error bounds, are
expected to give, but the specialization is not carried out here.

For \(1\le x\le2\), (P2) gives
\[
 |\ell_\lambda(\lambda x)|+
 \lambda^{-1}|\ell_\lambda'(\lambda x)|
 \ll\lambda^{1/2}\sqrt{d_4},
\tag{8.8n}
\]
which implies (8.8m) because \(x^{-3}\ge1/8\).

For \(2\le x\le c\), specialize Dunster's expansion to fixed order
\(j\in J_6\).  Its small parameter satisfies
\[
 \sigma_j^2=\frac{2j+1}{c}+O(c^{-2}).
\tag{8.8o}
\]
All indices are congruent modulo four, so the term \(-j\pi/2\) has the
same outgoing orientation.  Dunster's Liouville variable, denoted here
\(\zeta_D(\sigma,x)+J(\sigma)\) in the notation of \cite{DunsterPSWF},
has the expansion
\[
 \zeta_D(\sigma,x)+J(\sigma)
 =\sqrt{x^2-1}
 +\sigma^2\left(\frac\pi4-\frac12\operatorname {arcsec}x\right)
 +O\!\left(\sigma^4(1+x^{-1})\right).
\tag{8.8p}
\]
Consequently the radial expansion through order \(c^{-1}\) can be written
\[
 \widehat\psi_{j,\lambda}(\lambda x)
 =e_j\bigl(U_{0,\lambda}(x)+\mu_jU_{1,\lambda}(x)\bigr)
 +E_{j,\lambda}(x),
\tag{8.8q}
\]
where, by (P2), the two \(U\)'s are independent of \(j\), and the Olver
remainder, also after one derivative, is
\[
 |E_{j,\lambda}(x)|+
 \lambda^{-1}|\partial_tE_{j,\lambda}(x)|
 \ll |e_j|\bigl(x^{-3}+c^{-2}x^{-1}\bigr).
\tag{8.8r}
\]
Heuristically, (8.8p) makes the mode-dependent phase
\((2j+1)/(2x)+O_j(x^{-3}+c^{-1}x^{-1})\), while the radial amplitude is
\(x^{-1}(1+O_j(x^{-2}+c^{-1}))\).  Thus its constant and affine terms in
\(2j+1\) correspond to the two symbols in (8.8q).  These symbols
cancel because
\[
 \sum_jc_je_j=0,
 \qquad
 \sum_jc_j\mu_je_j=0
\tag{8.8s}
\]
by the endpoint constraints.  Since \(e_j^2\asymp\lambda d_j\) by
(P1), Cauchy--Schwarz over the six indices and
\(\sum_jd_j|c_j|^2\asymp d_4\sum_j|c_j|^2=d_4\) (by (8.8d)) give
\(\sum_j|c_je_j|\ll\lambda^{1/2}\sqrt{d_4}\); together with
\(c^{-2}x^{-1}\le x^{-3}\) for \(x\le c\), summing (8.8q)--(8.8r)
proves (8.8m) throughout this region.

For the range \(x\ge c\) we indicate where (8.8v) comes from.  If the outgoing solution is
written
\[
 y_j(t)=e^{2\pi i\lambda t}
 \sum_{m\ge0}a_{m,j}t^{-m-1}+\text{conjugate},
\]
then, with \(A_j=4\pi^2\lambda^4-\mu_j\), its exact recurrence is
\[
 \begin{aligned}
 2i(2\pi\lambda)m a_{m,j}
 ={}&[m(m-1)+A_j]a_{m-1,j}\\
 &+2i(2\pi\lambda)\lambda^2(m-1)a_{m-2,j}\\
 &-\lambda^2(m-2)(m-1)a_{m-3,j}.
 \end{aligned}
\tag{8.8t}
\]
The first two combined coefficients vanish by (8.8s), and
\[
 \sum_jc_ja_{2,j}
 =-\frac1{8(2\pi\lambda)^2}
   \sum_jc_j\mu_j^2a_{0,j}.
\tag{8.8u}
\]
A Volterra majorant for the remaining coefficients, using
\(|\mu_j|\ll\lambda^2\), is expected to give the bound assumed in (P2),
\[
 |\ell_\lambda(t)|+\lambda^{-1}|\ell_\lambda'(t)|
 \ll\lambda^{5/2}\sqrt{d_4}\,t^{-3},
 \qquad x\ge c,
\tag{8.8v}
\]
which is stronger than (8.8m).  The three regions cover
\([\lambda,\infty)\).
\end{proof}

\begin{theorem}[Even Rayleigh residual]
\label{thm:even-rayleigh-residual}
Assume (P1)--(P4).  For the endpoint-corrected moving vector of
Definition~\ref{def:endpoint-corrected-source},
\[
 \boxed{|\mathscr R_L^+|\ll\lambda^7d_4.}
\tag{8.8w}
\]
\end{theorem}

\begin{proof}
Define, for \(v\ge\lambda\),
\[
 H_\lambda(v)=v^{1/2}\sum_{n\ge1}\ell_\lambda(nv),
 \qquad D=v\partial_v.
\tag{8.8x}
\]
Lemma~\ref{lem:uniform-exterior-two-jet} and absolute summation imply
\[
 |H_\lambda(v)|\ll
 \lambda^{7/2}\sqrt{d_4}\,v^{-5/2},
 \qquad
 |DH_\lambda(v)|\ll
 \lambda^{9/2}\sqrt{d_4}\,v^{-3/2}.
\tag{8.8y}
\]
For \(|\Re w|<1/2\), put
\[
 A_\lambda(w)=\int_\lambda^\infty
 H_\lambda(v)v^{-w}\,\frac{dv}{v}.
\tag{8.8z}
\]
The direct estimate for bounded ordinate and the identity
\[
 wA_\lambda(w)
 =H_\lambda(\lambda)\lambda^{-w}
 +\int_\lambda^\infty
 DH_\lambda(v)v^{-w}\,\frac{dv}{v}
\tag{8.8aa}
\]
for large ordinate yield, uniformly in the strip,
\[
 |A_\lambda(w)|
 \ll\frac{\lambda^{7/2}\sqrt{d_4}}
          {1+|\Im w|}.
\tag{8.8ab}
\]
Since \(R_\lambda^+(x)=(h(x)+h(-x))/\sqrt2\) with
\(h(x)=H_\lambda(e^x)1_{[L/2,\infty)}(x)\), its Mellin coordinate is
\[
 b_\lambda(w)=
 \frac{A_\lambda(w)+A_\lambda(-w)}{\sqrt2}.
\tag{8.8ac}
\]
By (8.8y), \(R_\lambda^+\) and its derivative decay like
\(e^{-3|x|/2}\), so (P4) applies; as \(R_\lambda^+\) is real and even and
the zeros are symmetric under \(\rho\mapsto\bar\rho\), it gives
\[
 \QW(R_\lambda^+,R_\lambda^+)
 =\sum_\rho b_\lambda(\rho-\tfrac12)^2.
\tag{8.8ad}
\]
The local zero-count estimate \(N(T+1)-N(T)\ll\log(3+T)\), together
with (8.8ab), makes this series absolutely convergent and gives
\[
 |\QW(R_\lambda^+,R_\lambda^+)|
 \ll\lambda^7d_4.
\tag{8.8ae}
\]
Finally, the definition of \(\mathscr R_L^+\) and (8.8e) give
\[
 |\mathscr R_L^+|
 =\frac{|\QW(R_\lambda^+,R_\lambda^+)|}
        {\|V_\lambda^+\|^2}
 \ll\lambda^7d_4,
\]
as claimed.
\end{proof}

\begin{lemma}[Dyadic complex-frequency Bessel bound]
\label{lem:dyadic-complex-bessel}
Let \(I_L=[-L/2,L/2]\), \(\lambda=e^{L/2}\), and let coefficients indexed
by the nontrivial zeta divisor satisfy
\[
 |a_\rho|\le \frac{B}{1+|\gamma|},
 \qquad \rho=\beta+i\gamma .
\tag{8.8af}
\]
Then the symmetric partial sums of
\[
 h(x)=\sum_\rho a_\rho
 e^{-i\gamma x}e^{(\beta-1/2)x}
\tag{8.8ag}
\]
converge in \(L^2(I_L)\), and
\[
 \boxed{\|h\|_{L^2(I_L)}
 \ll B\sqrt{\lambda(1+L)}.}
\tag{8.8ah}
\]
The implied constant is absolute for the Riemann divisor.
\end{lemma}

\begin{proof}
It is enough to treat the positive ordinates, since the negative ordinates
are obtained by conjugation.  Separate a bounded initial block and put
\[
 \mathcal Z_k=\{\rho:2^k\le\gamma<2^{k+1}\},
 \qquad k\ge0.
\tag{8.8ai}
\]
For \(\rho,\rho'\in\mathcal Z_k\), direct integration gives the Gram
entry
\[
 \begin{aligned}
 G_{\rho,\rho'}
 &=\int_{-L/2}^{L/2}
 e^{(\beta+\beta'-1)x}e^{-i(\gamma-\gamma')x}\,dx,\\
 |G_{\rho,\rho'}|
 &\ll\lambda\min\!\left(L,
 \frac1{|\gamma-\gamma'|}\right),
 \end{aligned}
\tag{8.8aj}
\]
where the second expression is read as \(L\) at coincident ordinates.
Partition the ordinate axis into unit intervals.  The unconditional local
zero count
\[
 N(T+1)-N(T)\ll\log(3+T)
\tag{8.8ak}
\]
shows that the interval containing \(\gamma\), together with its two
neighbors, contributes \(O(\lambda L(1+k))\) to a Gram row.  The unit
intervals at distance \(m\ge2\) contribute
\(O(\lambda(1+k)/m)\).  Hence
\[
 \sup_{\rho\in\mathcal Z_k}
 \sum_{\rho'\in\mathcal Z_k}|G_{\rho,\rho'}|
 \ll\lambda(1+L+k)(1+k).
\tag{8.8al}
\]
Schur's test and Riemann--von Mangoldt now give
\[
 \begin{aligned}
 \left\|\sum_{\rho\in\mathcal Z_k}a_\rho
 e^{-i\gamma x}e^{(\beta-1/2)x}\right\|_2^2
 &\ll\lambda(1+L+k)(1+k)
       \sum_{\rho\in\mathcal Z_k}|a_\rho|^2,\\
 \sum_{\rho\in\mathcal Z_k}|a_\rho|^2
 &\ll B^2\frac{1+k}{2^k}.
 \end{aligned}
\tag{8.8am}
\]
After taking square roots, the triangle inequality over the dyadic blocks
yields
\[
 \|h\|_2
 \ll B\sqrt\lambda
 \sum_{k\ge0}
 \frac{(1+k)\sqrt{1+L+k}}{2^{k/2}}
 \ll B\sqrt{\lambda(1+L)}.
\tag{8.8an}
\]
The same summable majorant makes the symmetric partial sums Cauchy in
\(L^2(I_L)\), proving the lemma.
\end{proof}

\begin{theorem}[Cross residual of the endpoint-corrected source]
\label{thm:cross-residual-vanishing}
Assume (P1)--(P4).  For the endpoint-corrected moving vector of
Definition~\ref{def:endpoint-corrected-source},
\[
 \boxed{
 \|b_L^+\|
 \ll\lambda^4\sqrt{(1+L)d_4}.}
\tag{8.8ao}
\]
In particular, \(\|b_L^+\|\to0\).
\end{theorem}

\begin{proof}
Put
\[
 e_L^+=\frac{R_\lambda^+}{\|V_\lambda^+\|}.
\tag{8.8ap}
\]
Define the functional \(b_L^+\) on even \(g\in L^2(I_L)\) with
\(g\perp q_L^+\) by
\[
 \langle b_L^+,g\rangle=-\QW(e_L^+,g).
\tag{8.8aq}
\]
Thus \(\|b_L^+\|\) is the norm of this functional, as defined above.
If \(z_\rho=\rho-1/2\), the positive-parity Mellin coordinate is
\[
 \widehat {R_\lambda^+}(z_\rho)
 =b_\lambda(z_\rho),
 \qquad
 b_\lambda(w)=\frac{A_\lambda(w)+A_\lambda(-w)}{\sqrt2}.
\tag{8.8ar}
\]
Equations (8.8ab), (8.8e), \(0<\Re\rho<1\) and (P4) for the pairing
therefore give the complete-divisor representation
\[
 \QW(e_L^+,g)
 =\sum_\rho c_{\rho,\lambda}\widehat g(z_\rho),
 \qquad
 |c_{\rho,\lambda}|
 \ll\frac{B_\lambda}{1+|\gamma|},
 \quad
 B_\lambda=\lambda^{7/2}\sqrt{d_4}.
\tag{8.8as}
\]
Lemma~\ref{lem:dyadic-complex-bessel} shows that this functional has an
\(L^2(I_L)\) Riesz representative \(h_\lambda\) satisfying
\[
 \|h_\lambda\|_2
 \ll B_\lambda\sqrt{\lambda(1+L)}
 =\lambda^4\sqrt{(1+L)d_4}.
\tag{8.8at}
\]
Taking the supremum of (8.8aq) over the even unit ball orthogonal to
\(q_L^+\) proves (8.8ao).  Finally,
\[
 d_4\asymp\lambda^9e^{-4\pi\lambda^2}
\tag{8.8au}
\]
implies
\[
 \lambda^4\sqrt{(1+L)d_4}
 \asymp\lambda^{17/2}\sqrt{1+L}\,e^{-2\pi\lambda^2}
 \longrightarrow0.
\tag{8.8av}
\]
\end{proof}

Theorem~\ref{thm:cross-residual-vanishing} gives the cross-residual
bound used in Section~\ref{sec:cofinal-inertia}.  Together with
Theorem~\ref{thm:even-rayleigh-residual}, it gives, under (P1)--(P4),
\[
 |\mathscr R_L^+|+\|b_L^+\|\longrightarrow0.
\tag{8.8aw}
\]

\section{Cofinal prime--Gamma inertia matching}
\label{sec:cofinal-inertia}

The preceding moving-source analysis gives the two residual estimates under the standing assumptions (P1)--(P4).
The individual source components can exchange mass.  We therefore keep
all sources in one finite Hermitian form.

\subsection{Complete and finite joint heads}

On the centered real-even space
\[
 \mathcal H_0=\left\{q\in L^2_{\rm even}(\mu_K):\int q\,\dd\mu_K=0\right\},
\tag{9.1}
\]
define the polarized displacement atom
\[
\begin{aligned}
 \mathcal J_u(q,s)
 =\int_{\mathbb R}K(x)K(x-u)\,
 &\overline{q(x)-q(x-u)}\\
 &\times\bigl(s(x)-s(x-u)\bigr)\,\dd x .
\end{aligned}
\tag{9.2}
\]
The archimedean jump form is
\[
 \mathscr E_\Gamma(q,s)=\int_0^\infty g_\Gamma(u)\,\mathcal J_u(q,s)\,\dd u,
 \qquad
 g_\Gamma(u)=\frac{e^{u/2}}{e^{u}-e^{-u}} ,
\]
so that \(g_\Gamma(u)=\frac1{2u}+\frac14+O(u)\) as \(u\downarrow0\); relative to the
prime atoms below it carries the archimedean weight of the explicit
formula (1.1).  For \(X\ge2\), put
\[
\begin{aligned}
 \mathcal A_X(q,s)
 ={}&\mathscr E_\Gamma(q,s)
 +\sum_{p^k\le X}\frac{\Lambda(p^k)}{p^{k/2}}
       \mathcal J_{k\log p}(q,s)
 -\frac12\langle q,s\rangle_{\mu_K},\\
 \mathcal T_X(q,s)
 ={}&\sum_{p^k>X}\frac{\Lambda(p^k)}{p^{k/2}}
       \mathcal J_{k\log p}(q,s).
\end{aligned}
\tag{9.3}
\]
The complete form is
\[
 \mathcal A_\infty
 =\mathscr E_K-\frac12\langle\cdot,\cdot\rangle_{\mu_K},
\qquad
 \boxed{\mathcal A_X=\mathcal A_\infty-\mathcal T_X,\quad
        \mathcal T_X\succeq0.}
\tag{9.4}
\]
Here, and throughout this section, \(\mathscr E_K\) denotes the jump form
\(\mathscr E_\Gamma+\sum_{n}\frac{\Lambda(n)}{\sqrt n}\mathcal J_{\log n}\)
(its closure is constructed in Subsection~\ref{subsec:heat-core}).  Its
agreement on even \(r\in C_c^\infty(\mathbb R)\) with the jump energy of
Problem~\ref{prob:doob}, namely the ground-state identity
\[
 \mathscr E_\Gamma(r,r)+\sum_n\frac{\Lambda(n)}{\sqrt n}
 \mathcal J_{\log n}(r,r)
 =\QW(Kr,Kr)+\frac12\Var_{\mu_K}(r),
\]
is not proved in this paper; it is stated as a hypothesis wherever it is
used.
Thus the finite heads increase in Loewner order toward the complete form.

Let \(V_M\subset\mathcal H_0\) be a nested finite-dimensional mode space
in the complement of the complete radical.  The elementary choice is
generated by centered projections of
\[
 \chi_z(x)=\frac{\cos(zx)}{\cosh(x/2)}
\tag{9.5}
\]
for complete zero orbits \(z\) in a growing spectral rectangle, with all
multiplicity jets retained.  For a zero \(z\) of \(\XiR\) one has
\(\int\chi_z\,\dd\mu_K=2\int\cos(zx)K(x)\,\dd x=2\XiR(z)=0\), and
similarly for the multiplicity jets \(\partial_z^m\chi_z\) at a zero of
order \(>m\); so these modes are already centered.  Likewise, for the
functions \(r_j\) of (9.14) below, integration by parts gives
\(\langle\chi_z,r_j\rangle_{\mu_K}=2((-1)^jz^{2j}-4^{-j})\XiR(z)=0\).
Every element of
\(V_M\) is a finite combination of them, hence tends to zero at
\(\pm\infty\) because \(\lvert\operatorname{Im}z\rvert<\frac12\).
Let \(Z_M:\mathbb C^{d_M}\to V_M\) be the
synthesis map, and write
\[
 \mathbf N_M=Z_M^*Z_M,\qquad
 \mathbf H_{M,X}=Z_M^*\mathcal A_XZ_M .
\tag{9.6}
\]

\begin{theorem}[Cofinal filter-bank approximation]
\label{thm:cofinal-filter-bank}
There are constants \(c>0\) and \(C_M<\infty\) such that, for \(Y\ge X\),
\[
 \boxed{
 0\preceq
 \mathbf N_M^{-1/2}
 (\mathbf H_{M,\infty}-\mathbf H_{M,Y})
 \mathbf N_M^{-1/2}
 \preceq
 \mathbf N_M^{-1/2}
 (\mathbf H_{M,\infty}-\mathbf H_{M,X})
 \mathbf N_M^{-1/2}
 \preceq C_Me^{-cX}I.}
\tag{9.7}
\]
Consequently, for arbitrary \(\varepsilon_M\downarrow0\), the choice
\[
 X(M)\ge\frac1c
 \left(\log C_M+\log\frac1{\varepsilon_M}\right)
\tag{9.8}
\]
makes the omitted literal Euler bank at most \(\varepsilon_M\) in
normalized operator norm.
\end{theorem}

\begin{proof}
The first two inequalities follow from (9.3): each omitted prime-power
atom is a positive Gram form.  For fixed \(M\), every mode in (9.5) has
imaginary part strictly inside the critical strip.  The theta expansion of
\(K(x)K(x-u)\), with \(u=\log p^k\), is therefore bounded on \(V_M\) by a
fixed polynomial times \(\exp(-c p^k)\).  Summation over prime powers
larger than \(X\), followed by conjugation with \(\mathbf N_M^{-1/2}\),
gives the last inequality.  Solving \(C_Me^{-cX}\le\varepsilon_M\)
gives (9.8).
\end{proof}

Define the normalized bank \(S_{M,X}\) by
\[
 S_{M,X}^*S_{M,X}
 =\mathbf N_M^{-1/2}
  \left(\mathscr E_\Gamma+
  \sum_{p^k\le X}\frac{\Lambda(p^k)}{p^{k/2}}
       \mathcal J_{k\log p}\right)_{V_M}
 \mathbf N_M^{-1/2}.
\tag{9.9}
\]
Then the exact signed filter-bank identity is
\[
 \boxed{
 \mathbf G_{M,X}:=
 \mathbf N_M^{-1/2}\mathbf H_{M,X}\mathbf N_M^{-1/2}
 =S_{M,X}^*S_{M,X}-\frac12I.}
\tag{9.10}
\]
No scalar averaging occurs in (9.10): all modes and all retained sources
remain coupled in one matrix.

\subsection{The strict finite frontier}

Put
\[
 \delta_M=
 \lambda_{\min}\!\left(
 \mathbf N_M^{-1/2}\mathbf H_{M,\infty}
 \mathbf N_M^{-1/2}\right),
\qquad
 X_*(M)=\inf\{X:\mathbf H_{M,X}\succeq0\}.
\tag{9.11}
\]

\begin{theorem}[Exact finite-repair criterion]
\label{thm:exact-finite-repair}
For every fixed \(M\),
\[
 \boxed{X_*(M)<\infty\quad\Longleftrightarrow\quad\delta_M>0.}
\tag{9.12}
\]
If \(\delta_M>0\), then
\[
 X_*(M)\le
 \left\lceil\frac1c\log\frac{C_M}{\delta_M}\right\rceil .
\tag{9.13}
\]
If \(\delta_M=0\), every proper finite head has a negative direction and
its least eigenvalue increases to \(0\) without crossing.  If
\(\delta_M<0\), no finite head repairs the completed negative direction.
\end{theorem}

\begin{proof}
If \(\delta_M>0\), Theorem~\ref{thm:cofinal-filter-bank} gives
\[
 \mathbf H_{M,X}\succeq
 (\delta_M-C_Me^{-cX})\mathbf N_M,
\]
which proves (9.13).  Conversely, suppose
\(\mathbf H_{M,X}\succeq0\).  The omitted tail is positive definite on
\(V_M\): if it vanished on a nonzero \(q\), then the displacement atoms
of two omitted distinct primes would make \(q\) periodic with periods
\(\log p,\log r\).  Their ratio is irrational, so continuity makes \(q\)
constant; centeredness then makes \(q=0\), a contradiction.  Hence
\[
 \mathbf H_{M,\infty}
 =\mathbf H_{M,X}
  +(\mathbf H_{M,\infty}-\mathbf H_{M,X})\succ0.
\]
Finite dimensionality gives \(\delta_M>0\).  Evaluating a finite head on a
minimizer of (9.11) proves the last two assertions.
\end{proof}

\begin{remark}[What discreteness proves]
The finite spectra in (9.11) are discrete, but discreteness alone does not
imply (9.12)'s left side.  The scalar sequence \(-e^{-cX}\) is already a
discrete level which increases to \(0\) and never crosses it.  The
arithmetic theorem must therefore prove a strict completed gap or use the
vanishing compensated frontier below.
\end{remark}

\subsection{Radical shorting and the inertia staircase}

The word shorting is used in the Krein--Anderson--Trapp sense:
one eliminates an indefinite or constrained block through the extremal
shorted/compressed operator before reading the remaining Schur complement
\cite{Krein1947,Anderson1971,AndersonTrapp1975}.  The radical itself is
part of the Sonin/co-Poisson and Weil-distribution Hilbert-space
background cited in Remark~\ref{rem:attribution}; the explicit finite coordinate
\(r_j=K^{(2j)}/K-4^{-j}\) is the coordinate used here for the completed
radical.

Use a distinct symbol for the exact completed radical:
\[
 r_j=\frac{K^{(2j)}}K-4^{-j},
\qquad
 \mathfrak R_J=\operatorname{span}\{r_1,\dots,r_J\}.
\tag{9.14}
\]
Each \(r_j\) satisfies \(\int r_j\,\dd\mu_K=0\), since integrating by parts gives
\[
 \int\cosh(x/2)K^{(2j)}(x)\,\dd x=4^{-j}\int\cosh(x/2)K(x)\,\dd x .
\]
Moreover \(Kr_j=K^{(2j)}-4^{-j}K\) has Fourier transform
\(((-1)^jz^{2j}-4^{-j})\XiR(z)\), which vanishes at every zero of
\(\XiR\) and at \(z=\pm i/2\).  Assume the ground-state identity stated
after (9.4), and let \(f\) be even in \(C_c^\infty(\mathbb R)\).  The
explicit formula \(W(\phi)=\sum_\gamma\widehat\phi(\gamma)\), applied to
\(\phi=Kr_j*\widetilde{Kf}\), which is smooth and decays like
\(\exp(-\pi e^{2|x|})\) so that the formula applies, gives
\(\QW(Kr_j,Kf)=0\).  Polarizing the ground-state identity (extended to
\(r_j\) by the cutoff approximation of Subsection~\ref{subsec:heat-core},
under which both sides are continuous) and using \(\int r_j\,\dd\mu_K=0\)
gives \(\mathscr E_K(r_j,f)=\frac12\langle r_j,f\rangle_{\mu_K}\), i.e.
\[
 \mathcal A_\infty(r,f)=0
 \qquad(r\in\bigcup_J\mathfrak R_J,\ f\in C_{c,\rm even}^\infty(\mathbb R)),
\tag{9.15}
\]
and by continuity in the closed form norm for every \(f\) in the form
domain of Subsection~\ref{subsec:heat-core}.  Thus (9.15), and every
statement below that uses it, is conditional on the ground-state identity.
In bases of \(V_M\) and \(\mathfrak R_J\), write the omitted-tail blocks
as \(\mathbf T^{VV},\mathbf T^{VR},\mathbf T^{RR}\).  The radical block is
strictly positive:
\[
 \mathbf T^{RR}\succ0.
\tag{9.16}
\]
Indeed, vanishing would again give two incommensurable prime periods,
whereas the largest derivative in a nonzero finite combination of the
\(r_j\) has nonconstant theta asymptotics.

The finite-head matrix is therefore
\[
 \mathbf A_{M,J,X}=
 \begin{pmatrix}
  \mathbf H_M-\mathbf T^{VV}&-\mathbf T^{VR}\\
  -\mathbf T^{RV}&-\mathbf T^{RR}
 \end{pmatrix},
\qquad \mathbf H_M=\mathbf H_{M,\infty},
\tag{9.17}
\]
and its Schur complement is
\[
 \mathbf S_{M,J,X}
 =\mathbf H_M-
 \left(
 \mathbf T^{VV}
 -\mathbf T^{VR}(\mathbf T^{RR})^{-1}\mathbf T^{RV}
 \right).
\tag{9.18}
\]

\begin{theorem}[Finite inertia matching~\cite{Haynsworth1968,Gantmacher1959}]
\label{thm:finite-inertia-matching}
One has
\[
 \boxed{
 n_-(\mathbf A_{M,J,X})
 =J+n_-(\mathbf S_{M,J,X}).}
\tag{9.19}
\]
Consequently,
\[
 \boxed{
 n_-(\mathbf A_{M,J,X})=J
 \Longleftrightarrow
 \mathbf T^{VV}
 -\mathbf T^{VR}(\mathbf T^{RR})^{-1}\mathbf T^{RV}
 \preceq\mathbf H_M.}
\tag{9.20}
\]
\end{theorem}

\begin{proof}
Block Gaussian elimination is a congruence from (9.17) to
\(\mathbf S_{M,J,X}\oplus(-\mathbf T^{RR})\).  Equations
(9.16), Sylvester's law of inertia, and (9.18) give (9.19)--(9.20).
\end{proof}

Define the excess negative index and its frontier by
\[
 \kappa(M,J,X)=n_-(\mathbf A_{M,J,X})-J,
\qquad
 X_*(M,J)=\inf\{X:\kappa(M,J,X)=0\}.
\tag{9.21}
\]
Loewner monotonicity and interlacing give
\cite{Bhatia1997}
\[
\begin{array}{c|c}
\text{parameter increased}&\text{effect on }\kappa\\ \hline
X&\text{nonincreasing}\\
M&\text{nondecreasing}\\
J&\text{nonincreasing}.
\end{array}
\tag{9.22}
\]
Thus \(X_*(M+1,J)\ge X_*(M,J)\), whereas
\(X_*(M,J+1)\le X_*(M,J)\).  This is the exact mode--prime--radical
staircase: a head that closes a row with cutoff \(X\) may fail when the
next mode is admitted.

\subsection{Geometry of a large-prime channel}

The staircase has a literal signal-processing interpretation.  Recall
that \(K\) is normalized by \(\int K(x)e^{izx}\,\dd x=\XiR(z)\), so that
\[
 K(x)=2\sum_{n\ge1}\bigl(2\pi^2n^4e^{9x/2}-3\pi n^2e^{5x/2}\bigr)
 e^{-\pi n^2e^{2x}},
 \qquad
 K(x)=4\pi^2e^{9x/2}e^{-\pi e^{2x}}\bigl(1+O(e^{-2x})\bigr)
\]
as \(x\to+\infty\), and put \(C_\Xi=2\).

\begin{proposition}[Large-prime channel asymptotics]
For
\(\chi_z\) in (9.5), uniformly on fixed compact \(z,w\)-blocks with
\(\lvert\operatorname{Im}z\rvert,\lvert\operatorname{Im}w\rvert<\frac12\),
\[
 \boxed{
 \frac{\log p}{\sqrt p}\mathcal J_{\log p}(\chi_z,\chi_w)
 =\beta_p\,\overline{A_p(z)}A_p(w)+E_p(z,w),}
\tag{9.23}
\]
where
\[
\begin{aligned}
 A_p(z)
 &=2z\sin\!\left(\frac{z\log p}{2}\right)
 +\tanh\!\left(\frac{\log p}{4}\right)
  \cos\!\left(\frac{z\log p}{2}\right),\\
 \beta_p
 &=C_\Xi^2\pi^3(\log p)p^2e^{-2\pi p}
   (1+p^{-1/2})^{-2},
\end{aligned}
\tag{9.24}
\]
and
\[
 |E_p(z,w)|\le C_B\beta_pp^{b-1},
 \qquad b=\max(|\operatorname{Im}z|,|\operatorname{Im}w|)<\frac12.
\tag{9.25}
\]
\end{proposition}

\begin{proof}
Put \(a=\frac12\log p\) and \(x=a+t\).  Since \(K\) and \(\chi_z\) are
even, \(K(x)K(x-\log p)=K(a+t)K(a-t)\) and
\(\chi_z(x)-\chi_z(x-\log p)=\chi_z(a+t)-\chi_z(a-t)\).  For
\(\lvert t\rvert\le\frac14\log p\) the asymptotics of \(K\) give
\[
 K(a+t)K(a-t)=16\pi^4p^{9/2}e^{-2\pi p\cosh 2t}
 \bigl(1+O(p^{-1}\cosh 2t)\bigr),
\]
because \(e^{9(a+t)/2}e^{9(a-t)/2}=p^{9/2}\) and
\(e^{2(a+t)}+e^{2(a-t)}=2p\cosh2t\); the terms \(n\ge2\) of the theta
series are smaller by a factor \(O(e^{-3\pi\sqrt p})\) on this range.  On
\(\lvert t\rvert\ge p^{-1/2+\epsilon}\) (in particular on
\(\lvert t\rvert>\frac14\log p\), where \(e^{2|t|}\ge p^{1/2}\)) the
factor \(e^{-2\pi p(\cosh2t-1)}\le e^{-4\pi pt^2}\) makes the
contribution smaller than any power of \(p\) relative to the main term,
since \(\chi_z\) grows at most like \(e^{b|x|}\).  The weight is even
in \(t\), and
\(\chi_z(a+t)-\chi_z(a-t)=2t\chi_z'(a)+O(p^{b/2-1/4}\lvert t\rvert^3)\)
with
\[
 \chi_z'(a)=-\frac{z\sin(za)+\frac12\tanh(a/2)\cos(za)}{\cosh(a/2)}
 =-\frac{A_p(z)}{p^{1/4}(1+p^{-1/2})},
\]
because \(\cosh(a/2)=\frac12p^{1/4}(1+p^{-1/2})\).  Using
\(e^{-2\pi p\cosh2t}=e^{-2\pi p}e^{-4\pi pt^2}(1+O(pt^4))\) and
\(\int_{\mathbb R}4t^2e^{-4\pi pt^2}\,\dd t=\frac1{4\pi p^{3/2}}\), the
main term is
\[
 16\pi^4p^{9/2}e^{-2\pi p}\,\frac{\overline{A_p(z)}A_p(w)}
 {p^{1/2}(1+p^{-1/2})^2}\cdot\frac1{4\pi p^{3/2}}
 =4\pi^3p^{5/2}e^{-2\pi p}\frac{\overline{A_p(z)}A_p(w)}{(1+p^{-1/2})^2},
\]
and multiplication by \((\log p)/\sqrt p\) gives
\(\beta_p\overline{A_p(z)}A_p(w)\) with \(C_\Xi^2=4\).  On the
effective range \(t=O(p^{-1/2+\epsilon})\), the cubic Taylor remainder
(relative size \(t^2\)), the quartic term of \(\cosh2t\) (relative size
\(pt^4\)) and the correction factor of \(K\) (relative size \(p^{-1}\))
are each \(O(p^{-1+4\epsilon})\) relative to a quantity bounded by
\(\beta_pp^{b}\); here \(|\chi_z'(a)|,|\chi_z'''(a)|\ll p^{b/2-1/4}\) on
the compact block.  Taking the Gaussian moments exactly instead of
estimating them removes the \(\epsilon\), and gives (9.25).
\end{proof}

Numerically, the ratio of the left side of (9.23) to
\(\beta_p\overline{A_p(z)}A_p(w)\) is approximately \(0.867\),
\(0.953\), \(0.985\) at \(p=11,29,101\) for \(z=w=3\), consistent with
(9.23)--(9.25).
Heuristically, each large prime contributes approximately one derivative
direction; this is not used below.

\subsection{Finite observability and the exact gain budget}

The distinction between observability and gain can be made exact.  Every
single literal displacement observes every finite elementary mode space.

\begin{theorem}[One-channel finite observability]
\label{thm:one-channel-observability}
For every \(u>0\) and every nonzero \(q\in V_M\),
\[
 \boxed{\mathcal J_u(q,q)>0.}
\tag{9.25a}
\]
Consequently the atom Gram matrix is positive definite on \(V_M\).  Its
normalized least eigenvalue is bounded away from zero uniformly when
\(u\) ranges over a fixed compact subinterval of \((0,\infty)\).
\end{theorem}

\begin{proof}
The real theta kernel is strictly positive.  Therefore
\(\mathcal J_u(q,q)=0\) would imply
\(q(x)=q(x-u)\) for almost every \(x\), and analyticity extends the equality
to the real line.  As noted after (9.5), every element of \(V_M\) is a
finite combination of the centered modes \(\chi_z\) and their
multiplicity jets, and tends to zero at both ends because
\(\lvert\operatorname {Im}z\rvert<1/2\).  A periodic function with this
decay is zero.  This proves strict positivity.  Continuity in \(u\) and
compactness of the normalized unit sphere give the uniform assertion.
\end{proof}

Write the normalized Gamma deficit and prime bank as
\[
 \mathbf B_M=\frac12I-
 \mathbf N_M^{-1/2}(\mathscr E_\Gamma)_{V_M}\mathbf N_M^{-1/2},
 \qquad
 \mathbf P_{M,X}=\sum_{p^k\le X}\frac{\log p}{p^{k/2}}
 \mathbf N_M^{-1/2}(\mathcal J_{k\log p})_{V_M}\mathbf N_M^{-1/2}.
\tag{9.25b}
\]
Then \(\mathbf G_{M,X}=\mathbf P_{M,X}-\mathbf B_M\), and congruence by
\(\mathbf P_{M,X}^{-1/2}\) gives
\[
 \boxed{
 \mathbf G_{M,X}\succeq0
 \Longleftrightarrow
 \vartheta_{M,X}:=
 \lambda_{\max}(\mathbf P_{M,X}^{-1/2}
 \mathbf B_M\mathbf P_{M,X}^{-1/2})\le1.}
\tag{9.25c}
\]
Whenever a new prime-power atom is added, Theorem
\ref{thm:one-channel-observability} makes \(\vartheta_{M,X}\) strictly
decrease on the positive-deficit sector.  Norm convergence gives
\[
 \boxed{X_*(M)<\infty\quad\Longleftrightarrow\quad
 \vartheta_{M,\infty}<1.}
\tag{9.25d}
\]
Thus finiteness of \(X_*(M)\) is equivalent to the strict inequality in
(9.25d).  The strict inequality in (9.25d) is the gain question; it
cannot be replaced by the sum of the individual least atom eigenvalues,
because that replacement discards the cross-prime alignment.

\subsection{Scalar innovations of the completed staircase}

There is an exact one-scalar recursion for the remaining matrix sign.
Choose a basis \(\phi_1,\phi_2,\ldots\) adapted to the nested spaces and
write
\[
 \mathbf H_M=
 \begin{pmatrix}\mathbf H_{M-1}&c_M\\c_M^*&h_M\end{pmatrix}.
\tag{9.25e}
\]
Assuming \(\mathbf H_{M-1}\succ0\), set
\[
 a_M=\mathbf H_{M-1}^{-1}c_M,
 \qquad
 q_M^*=\phi_M-\sum_{j<M}(a_M)_j\phi_j.
\tag{9.25f}
\]

\begin{theorem}[Completed scalar innovation]
\label{thm:scalar-innovation}
The vector \(q_M^*\) is orthogonal to \(V_{M-1}\) in the complete signed
form and
\[
\boxed{
 \begin{aligned}
 \sigma_M
 &:={\mathcal A}_\infty(q_M^*,q_M^*)
 =h_M-c_M^*\mathbf H_{M-1}^{-1}c_M
 =\frac{\det\mathbf H_M}{\det\mathbf H_{M-1}},\\
 \mathbf H_M\succ0
 &\Longleftrightarrow
 \mathbf H_{M-1}\succ0\ \text{and}\ \sigma_M>0.
 \end{aligned}}
\tag{9.25g}
\]
Moreover the scalar retains the complete physical source decomposition
\[
 \boxed{
 \sigma_M=\mathscr E_\Gamma(q_M^*)
 +\sum_{p^k}\frac{\log p}{p^{k/2}}
       \mathcal J_{k\log p}(q_M^*,q_M^*)
 -\frac12\lVert q_M^*\rVert_{\mu_K}^2.}
\tag{9.25h}
\]
\end{theorem}

\begin{proof}
Block multiplication gives
\({\mathcal A}_\infty(v,q_M^*)=0\) for \(v\in V_{M-1}\) and the first
line of (9.25g).  Congruence by triangular block elimination reduces
\(\mathbf H_M\) to \(\mathbf H_{M-1}\oplus[\sigma_M]\); Sylvester's law
and the block determinant formula give the second line.  Substitution in
(9.3) gives (9.25h).
\end{proof}

Before radical anti-shorting, the strict frontier is finite in every row
precisely when all literal innovations (9.25h) are positive.  This scalar
criterion is an exact unshorted precursor: it keeps Gamma, the threshold,
all ordinary von Mangoldt weights and their directional alignment in one
number for each newly admitted mode.  The active shortened target is
identified below by \(\mathfrak g_{M,J,X}>\delta^{\flat}_{M,J,X}\).

The moving-filter interpretation also has an exact update law.  Suppose
the preceding-mode block \(A\) is positive before the atom
\(n=p^k\) is inserted, write \(w_n=\Lambda(n)/\sqrt n\), let
\(U_n=D_n|_{V_{M-1}}\), and put
\[
 D_nq(x)=\sqrt{K(x)K(x-\log n)}
          \{q(x)-q(x-\log n)\}.
\tag{9.25i}
\]
If \(q_{M,X}^*\) is the adaptive finite-head regression residual and
\(r_n=D_nq_{M,X}^*\), direct Schur minimization gives
\[
 \boxed{
 \sigma_{M,+}-\sigma_{M,-}
 =w_n\left\langle r_n,
 (I+w_nU_nA^{-1}U_n^*)^{-1}r_n\right\rangle>0.}
\tag{9.25j}
\]
Indeed, after writing an old-mode correction as \(y\), the new pivot is
the minimum of
\[
 \sigma_{M,-}+\langle y,Ay\rangle
 +w_n\lVert r_n+U_ny\rVert^2;
\]
the normal equation and the Woodbury identity give (9.25j).  Iterating
over the ordered prime powers yields
\[
 \sigma_{M,X_L}=\sigma_{M,X_0}
 +\sum_{X_0<n_\ell\le X_L}\Delta_{M,\ell},
 \qquad \Delta_{M,\ell}>0.
\tag{9.25k}
\]
Thus the row closes at the first index at which this exact cumulative
innovation exceeds the initial deficit.  The series converges to
\(\sigma_M\); it cannot be replaced by a divergent ``infinitely many
channels'' argument, because theta localization makes the large-prime
gain summable.

Several primes can be kept jointly.  For a finite block \(\mathcal B\),
put
\[
 \mathcal D_{\mathcal B}q
 =\bigoplus_{n\in\mathcal B}\sqrt{\frac{\Lambda(n)}{\sqrt n}}D_nq,
 \quad U_{\mathcal B}=\mathcal D_{\mathcal B}|_{V_{M-1}},
 \quad r_{\mathcal B}=\mathcal D_{\mathcal B}q_{M,X}^*.
\tag{9.25l}
\]
The same minimization proves the block update
\[
 \Delta_{\mathcal B}
 =\langle r_{\mathcal B},
 (I+U_{\mathcal B}A^{-1}U_{\mathcal B}^*)^{-1}
 r_{\mathcal B}\rangle
 \ge\frac{\lVert r_{\mathcal B}\rVert^2}
 {1+\lVert U_{\mathcal B}\rVert^2/\lambda_{\min}(A)}.
\tag{9.25m}
\]
Unlike a sum of atomwise floors, (9.25m) retains complementarity between
prime sensors.  It gives a finite interval-certifiable crossing criterion:
partition the unused prime powers into finitely many blocks and verify that
the sum of their lower bounds exceeds the current negative pivot.

\begin{remark}[Positive weights alone do not force a crossing]
Positive weights, strict observability and norm convergence alone do not
force the strict inequality in (9.25d).  In a three-dimensional model,
let \(G\) be \(\frac18\) times the path Laplacian on three vertices
(centered eigenvalues \(\frac18,\frac38\)), let \(\Pi\) project onto the
centered plane, and choose positive numbers \(a_n\), indexed by prime
powers, with \(\sum a_n=\frac18\); each channel \(a_n\Pi\) is strictly
positive on the centered plane.  The completed energies of
\(G+\sum_na_n\Pi\) are
\[
 0,\qquad\frac14,\qquad\frac12,
\tag{9.25n}
\]
so the analogue of \(\mathcal A_\infty\), the energy minus
\(\frac12\), has the eigenvalue \(-\frac14\) on the centered plane.  Any
proof of (9.25d) must therefore use the specific theta geometry of
(9.25i).
\end{remark}

\subsection{Local detection and cumulative adaptive gain}

The strict observability statement can be localized at the midpoint of
one prime-power displacement.  For \(u,\eta>0\), put \(t=u/2\) and
define
\[
\begin{aligned}
 \mathcal J_{u,\eta}(q,s)
 :=\int_{-\eta}^{\eta}&K(t+y)K(t-y)
 \overline{\{q(t+y)-q(t-y)\}}\\
 &\times\{s(t+y)-s(t-y)\}\,\dd y .
\end{aligned}
\tag{9.25o}
\]
Let \(\mathbf C_M^{\rm loc}(u,\eta)\) be its Gram matrix in
\((\phi_1,\ldots,\phi_M)\).

\begin{proposition}[Local Christoffel detection~\cite{Bhatia1997}]
\label{prop:local-christoffel}
For every \(u,\eta>0\),
\[
 \mathbf C_M^{\rm loc}(u,\eta)\succ0,
 \qquad
 \boxed{\kappa_M(u,\eta)
 =\frac{\det\mathbf C_M^{\rm loc}(u,\eta)}
        {\det\mathbf C_{M-1}^{\rm loc}(u,\eta)}>0.}
\tag{9.25p}
\]
Every affine residual \(q=\phi_M+v\), \(v\in V_{M-1}\), satisfies
\(\mathcal J_{u,\eta}(q,q)\ge\kappa_M(u,\eta)\).  Hence an omitted
prime power \(n=p^k\) gives the adaptive lower bound
\[
 \Delta_n\ge\frac{\Lambda(n)}{\sqrt n}
 \kappa_M(\log n,\eta).
\tag{9.25q}
\]
\end{proposition}

\begin{proof}
If the local integral vanishes, strict positivity of \(K\) gives
\(q(t+y)=q(t-y)\) on an interval.  Analytic continuation makes \(q\)
invariant under reflection about \(t\).  Since \(q\) is even, composition
with reflection about zero makes it \(u\)-periodic.  By the remark after
(9.5), every vector of \(V_M\) tends to zero at \(+\infty\); a
\(u\)-periodic function with this property vanishes, so \(q=0\).  The local Gram is positive definite.
Minimizing its quadratic polynomial over the coefficient of
\(V_{M-1}\) gives the last Schur complement, which is the determinant
ratio in (9.25p).  Restricting the literal atom to the midpoint aperture
and dropping the nonnegative preceding-mode penalty gives (9.25q).
\end{proof}

The atomwise statement has an exact cumulative version.  Partition the
omitted prime powers into consecutive finite blocks
\(\mathcal B_1,\mathcal B_2,\ldots\).  At stage \(\ell\), let
\(A_{\ell-1}\succ0\) be the preceding-mode block, let \(U_\ell\) be the
old-mode feature of \(\mathcal B_\ell\), and let \(\rho_\ell\) be its
adaptive new-mode response.  Set
\[
 B_\ell=U_\ell A_{\ell-1}^{-1}U_\ell^*,\qquad
 R_\ell=\|\rho_\ell\|^2,
 \qquad
 C_\ell=\|A_{\ell-1}^{-1/2}U_\ell^*\rho_\ell\|^2.
\tag{9.25r}
\]

\begin{proposition}[Cumulative adaptive block gain]
\label{prop:cumulative-gain}
If the initial pivot is \(-\delta\), then after \(L\) blocks
\[
\boxed{
 \sigma_L=-\delta+\sum_{\ell=1}^L g_\ell,
 \qquad
 g_\ell=\langle\rho_\ell,(I+B_\ell)^{-1}\rho_\ell\rangle
 \ge\frac{R_\ell^2}{R_\ell+C_\ell}.}
\tag{9.25s}
\]
Consequently a finite prefix crosses if and only if the completed exact
gain satisfies \(\sum_{\ell\ge1}g_\ell>\delta\).  The sufficient
two-moment test is
\[
 \sum_{\ell\ge1}\frac{R_\ell^2}{R_\ell+C_\ell}>\delta.
\tag{9.25t}
\]
\end{proposition}

\begin{proof}
The block Woodbury identity gives
\(\sigma_\ell-\sigma_{\ell-1}=g_\ell\), so the first identity telescopes.
If \(\nu_\ell\) is the spectral measure of \(B_\ell\) at
\(\rho_\ell\), then
\[
 g_\ell=\int\frac{\dd\nu_\ell(s)}{1+s},\quad
 R_\ell=\int\dd\nu_\ell(s),\quad
 R_\ell+C_\ell=\int(1+s)\dd\nu_\ell(s).
\]
Cauchy--Schwarz gives \(R_\ell^2\le
g_\ell(R_\ell+C_\ell)\).  Monotonicity of the partial sums proves the
finite-selection assertion.
\end{proof}

Neither strict local detection nor (9.25s) by itself forces a crossing.
A nonzero completed radical has negative pivot equal to the omitted tail
at every proper cutoff and acquires exactly that amount only in the
limit.  The radical must therefore be removed before its tail response is
compared with the source deficit.

\subsection{Radical-conditioned tail gain}

Let \(\Phi:\C^{M-1}\to V_{M-1}\) and
\(\Psi:\C^J\to\mathfrak R_J\) be coordinate maps.  Realize the omitted
Euler bank by a Hilbert-space feature \(\mathcal V_X\) satisfying
\[
 \|\mathcal V_Xq\|^2=\mathcal T_X(q,q),
 \qquad
 U=\mathcal V_X\Phi,\quad v=\mathcal V_X\phi_M,
 \quad W_J=\mathcal V_X\Psi.
\tag{9.25u}
\]
Strict tail positivity on the finite radical space gives
\(W_J^*W_J\succ0\).  Put
\[
 \Pi_{J,X}=I-W_J(W_J^*W_J)^{-1}W_J^*.
\tag{9.25v}
\]
Thus \(\Pi_{J,X}\) removes precisely the tail features synthesized by a
finite exact-radical correction.

Write the completed mode matrix on
\(V_{M-1}\oplus\operatorname{span}\{\phi_M\}\) as \(H_\infty\).
Maximizing the finite-head form in the radical coordinate gives the exact
anti-short
\[
 \boxed{
 \widehat H_{M,J,X}=H_\infty-
 \begin{pmatrix}U^*\\v^*\end{pmatrix}
 \Pi_{J,X}\begin{pmatrix}U&v\end{pmatrix}
 =\begin{pmatrix}\widehat A&\widehat c\\
                  \widehat c^*&\widehat h\end{pmatrix}.}
\tag{9.25w}
\]
Assume \(\widehat A\succ0\), and define
\[
 \widehat a=\widehat A^{-1}\widehat c,\qquad
 q_{M,J,X}^*=\phi_M-\Phi\widehat a,\qquad
 \widehat\sigma_{M,J,X}
 =\widehat h-\widehat c^*\widehat A^{-1}\widehat c
 =:-\delta^{\flat}_{M,J,X}.
\tag{9.25x}
\]
The row is called negative if \(\delta^{\flat}_{M,J,X}>0\).  (The
symbol \(\delta^\flat\) is distinct from \(\delta_M\) of (9.11), which
has the opposite sign convention.)  The identities (9.25aa)--(9.25al)
below hold for either sign of \(\delta^{\flat}_{M,J,X}\).
The stationary radical coefficient and full saddle residual are
\[
 \zeta_{M,J,X}^*=-(W_J^*W_J)^{-1}W_J^*\mathcal V_Xq_{M,J,X}^*,
 \qquad
 \widetilde q_{M,J,X}^*=q_{M,J,X}^*+\Psi\zeta_{M,J,X}^*.
\tag{9.25y}
\]
In particular,
\[
 \mathcal V_X\widetilde q_{M,J,X}^*
 =\Pi_{J,X}\mathcal V_Xq_{M,J,X}^*.
\tag{9.25z}
\]

\begin{theorem}[Exact radical-conditioned adaptive gain]
\label{thm:radical-conditioned-gain}
For a row with \(\widehat A\succ0\), set
\[
 z_{M,J,X}=\Pi_{J,X}\mathcal V_Xq_{M,J,X}^*,
 \qquad \overline U=\Pi_{J,X}U.
\]
Then the completed last pivot is
\[
 \boxed{
 \sigma_{M,\infty}=-\delta^{\flat}_{M,J,X}+\mathfrak g_{M,J,X},
 \qquad
 \mathfrak g_{M,J,X}
 =\left\langle z_{M,J,X},
 (I+\overline U\widehat A^{-1}\overline U^*)^{-1}
 z_{M,J,X}\right\rangle.}
\tag{9.25aa}
\]
Equivalently,
\[
 \boxed{
 \mathfrak g_{M,J,X}
 =\sup_{\omega\ne0}
 \frac{|\langle\omega,z_{M,J,X}\rangle|^2}
 {\|\omega\|^2+
  \|\widehat A^{-1/2}\overline U^*\omega\|^2}.}
\tag{9.25ab}
\]
Finite-support tail combiners are sufficient in the strict inequality.
Thus
\[
 \boxed{\sigma_{M,\infty}>0
 \quad\Longleftrightarrow\quad
 \mathfrak g_{M,J,X}>\delta^{\flat}_{M,J,X}.}
\tag{9.25ac}
\]
\end{theorem}

\begin{proof}
The radical block in the joint finite-head matrix is
\(-W_J^*W_J\).  Its Schur complement is (9.25w), while its normal
equation gives (9.25y)--(9.25z).  Adding the omitted bank back to
\(\widehat H_{M,J,X}\) is therefore the positive update with old feature
\(\overline U\) and adaptive response \(z_{M,J,X}\).  The block Woodbury
identity gives (9.25aa).  The Rayleigh dual formula for the positive
operator \(I+\overline U\widehat A^{-1}\overline U^*\) gives (9.25ab).
Finite-support vectors are dense in the tail direct sum, and the two
quadratic functionals in (9.25ab) are continuous; a strict witness may
therefore be approximated by a finite-support witness.  Finally
(9.25ac) follows from (9.25aa).
\end{proof}

The formula also has the joint variational form
\[
 \boxed{
 \mathfrak g_{M,J,X}
 =\inf_{\substack{d\in V_{M-1}\\r\in\mathfrak R_J}}
 \left\{\widehat{\mathcal A}(d,d)
 +\|\mathcal V_X(q_{M,J,X}^*+d+r)\|^2\right\},}
\tag{9.25ad}
\]
where \(\widehat{\mathcal A}\) is the positive old-mode form with matrix
\(\widehat A\).  This identity keeps one common radical correction and
one common old-mode correction across all prime powers.

\subsection{Radical-conditioned Christoffel and finite determinant
exhaustion}

On the elementary analytic space \(V_M\oplus\mathfrak R_J\), every
literal midpoint form (9.25o) is positive definite.  Indeed, a null
vector is periodic by the reflection argument.  Its zero-mode component
decays, whereas a nonzero largest radical derivative has a nonconstant
leading polynomial in \(e^{2x}\); their sum cannot be periodic.  Hence,
if \(\mathbf C_{M,J}^{\rm loc}(u,\eta)\) is the local Gram in a basis of
\(V_M\oplus\mathfrak R_J\),
\[
 \boxed{
 \kappa_{M,J}(u,\eta)
 =\frac{\det\mathbf C_{M,J}^{\rm loc}(u,\eta)}
        {\det\mathbf C_{M-1,J}^{\rm loc}(u,\eta)}>0.}
\tag{9.25ae}
\]
Consequently every finite omitted block \(\mathcal B\) gives the valid
but not necessarily sharp certificate
\[
 \mathfrak g_{M,J,X}\ge
 \sum_{n\in\mathcal B}\frac{\Lambda(n)}{\sqrt n}
 \kappa_{M,J}(\log n,\eta_n).
\tag{9.25af}
\]
The separate minima in (9.25af) can lose the joint alignment.  The
following augmented determinant removes that loss.

For \(Y>X\), define
\[
\begin{aligned}
 \mathfrak C_{M,J;X,Y}:=
 \inf_{\substack{d\in V_{M-1}\\r\in\mathfrak R_J}}
 \bigg\{&\widehat{\mathcal A}(d,d)\\
 &+\sum_{\substack{n=p^k\\X<n\le Y}}
 \frac{\Lambda(n)}{\sqrt n}
 \mathcal J_{\log n}(q_{M,J,X}^*+d+r,q_{M,J,X}^*+d+r)
 \bigg\}.
\end{aligned}
\tag{9.25ag}
\]
In the ordered affine basis
\((V_{M-1},\mathfrak R_J;q_{M,J,X}^*)\), let
\(\mathbb G_n\) be the Gram matrix of \(\mathcal J_{\log n}\), and put
\[
 \mathbb C_{M,J;X,Y}
 =\operatorname{diag}(\widehat A,0_J,0)
 +\sum_{\substack{n=p^k\\X<n\le Y}}
 \frac{\Lambda(n)}{\sqrt n}\mathbb G_n.
\tag{9.25ah}
\]
Let \(\mathbb C_{M,J;X,Y}^{-}\) be the principal block obtained by
deleting the last row and column.

\begin{theorem}[Augmented finite determinant exhaustion]
\label{thm:augmented-exhaustion}
For a row with \(\widehat A\succ0\) and every \(Y\) containing at least one omitted prime power,
\[
 \boxed{
 \mathfrak C_{M,J;X,Y}
 =\frac{\det\mathbb C_{M,J;X,Y}}
        {\det\mathbb C_{M,J;X,Y}^{-}}>0.}
\tag{9.25ai}
\]
Moreover,
\[
 \boxed{
 0<\mathfrak C_{M,J;X,Y}<\mathfrak g_{M,J,X},\qquad
 \mathfrak C_{M,J;X,Y}\nearrow\mathfrak g_{M,J,X},}
\tag{9.25aj}
\]
and, on each fixed finite row, there are \(C,c>0\) and a polynomial
\(Q\) such that
\[
 0\le\mathfrak g_{M,J,X}-\mathfrak C_{M,J;X,Y}
 \le C Q(Y)e^{-cY}.
\tag{9.25ak}
\]
Consequently
\[
 \boxed{
 \sigma_{M,\infty}>0
 \quad\Longleftrightarrow\quad
 \mathfrak C_{M,J;X,Y}>\delta^{\flat}_{M,J,X}
 \text{ for some finite }Y.}
\tag{9.25al}
\]
\end{theorem}

\begin{proof}
A single displacement Gram is positive definite on
\(V_M\oplus\mathfrak R_J\) by the periodicity/asymptotic argument above.
The positive ordinary weights therefore make both the matrix in
(9.25ah) and its preceding principal block positive definite.  Its last
Schur complement is exactly the minimization (9.25ag), proving
(9.25ai).

The objectives in (9.25ag) increase pointwise with \(Y\) and are bounded
above by (9.25ad).  One fixed nonempty head is coercive on the
finite-dimensional correction space, so all finite minimizers lie in a
common compact set.  Passing to a convergent subsequence and then using
monotone convergence of the nonnegative prime-power series proves that
the limit is (9.25ad).  If \((d_\infty,r_\infty)\) minimizes the complete
objective, then \(q_{M,J,X}^*+d_\infty+r_\infty\ne0\); every omitted
literal atom is strictly positive on it.  Every proper cutoff therefore
omits a positive term, proving the strict inequality in (9.25aj).
The fixed-space theta-overlap estimate bounds the remaining Gram tail by
\(CQ(Y)e^{-cY}\); compactness of the minimizers gives (9.25ak).
Finally combine (9.25aa) and (9.25aj): a positive limiting margin is
eventually larger than the vanishing tail loss, while a finite crossing
implies \(\mathfrak g_{M,J,X}>\delta^{\flat}_{M,J,X}\).  This proves (9.25al).
\end{proof}

The finite source equation eliminates the abstract deficit.  Namely,
\[
\boxed{
\begin{aligned}
 \delta^{\flat}_{M,J,X}={}&\frac12\|\widetilde q_{M,J,X}^*\|_{\mu_K}^2
 -\mathscr E_\Gamma(\widetilde q_{M,J,X}^*)\\
 &-\sum_{p^k\le X}\frac{\log p}{p^{k/2}}
 \mathcal J_{k\log p}(\widetilde q_{M,J,X}^*,
                       \widetilde q_{M,J,X}^*).
\end{aligned}}
\tag{9.25am}
\]
There is an exact cancellation which rules out a raw-response shortcut.
Put
\[
 R_{M,J,X}=\|z_{M,J,X}\|^2,
 \qquad
 \mathfrak l_{M,J,X}=R_{M,J,X}-\mathfrak g_{M,J,X}\ge0.
\]
Radical invariance of the complete form, (9.25z), and
\(\mathcal A_X=\mathcal A_\infty-\mathcal T_X\) give
\[
 \boxed{
 R_{M,J,X}-\delta^{\flat}_{M,J,X}
 =\mathcal A_\infty(q_{M,J,X}^*,q_{M,J,X}^*),
 \qquad
 \mathfrak g_{M,J,X}-\delta^{\flat}_{M,J,X}
 =\mathcal A_\infty(q_{M,J,X}^*,q_{M,J,X}^*)
  -\mathfrak l_{M,J,X}.}
\tag{9.25an}
\]
Thus even the prerequisite \(R_{M,J,X}>\delta^{\flat}_{M,J,X}\) is already
positivity of the completed form on the current residual.  Bounds which
first dominate the raw response and then pay a positive regression penalty
are strictly stronger than the required statement.  The augmented
determinant is sharp because it retains \(\mathfrak l_{M,J,X}\) exactly.

The same assertion admits a source-balanced bordered-determinant form.
Write
\[
 \tau_{d+1}(Y)=\det\mathbb C_{M,J;X,Y},\qquad
 \tau_d(Y)=\det\mathbb C_{M,J;X,Y}^{-},
\]
let \(e_*\) denote the last coordinate vector, and set
\[
 \mathbb P_{M,J;X,Y}
 =\mathbb C_{M,J;X,Y}
  -\delta^{\flat}_{M,J,X}e_*e_*^*.
\]
Taking the Schur complement of the preceding principal block gives the
exact identity
\[
 \boxed{
 \det\mathbb P_{M,J;X,Y}
 =\tau_d(Y)\{
   \mathfrak C_{M,J;X,Y}-\delta^{\flat}_{M,J,X}\}
 =\tau_{d+1}(Y)-\delta^{\flat}_{M,J,X}\tau_d(Y).}
\]
Since \(\tau_d(Y)>0\), finite crossing is therefore precisely the strict
minor domination
\[
 \boxed{\tau_{d+1}(Y)>\delta^{\flat}_{M,J,X}\tau_d(Y).}
\]
This recoding introduces neither a localization loss nor a separate
effective-resistance estimate.

Thus (9.25al) is the completely finite inequality
\[
\begin{aligned}
 \frac{\det\mathbb C_{M,J;X,Y}}
      {\det\mathbb C_{M,J;X,Y}^{-}}
 &+\mathscr E_\Gamma(\widetilde q_{M,J,X}^*)\\
 &+\sum_{p^k\le X}\frac{\log p}{p^{k/2}}
 \mathcal J_{k\log p}(\widetilde q_{M,J,X}^*,
                       \widetilde q_{M,J,X}^*)
 >\frac12\|\widetilde q_{M,J,X}^*\|_{\mu_K}^2.
\end{aligned}
\tag{9.25ao}
\]
It retains Gamma, the scalar threshold, the exact radical correction, the
old-mode adaptation, every retained ordinary prime power, and one finite
omitted block in a single determinant.

\subsection{Signed amplitude transfer and adaptive Krylov recovery}
\label{sec:signed-krylov}

Fix one row \((M,J,X,Y)\), put \(d=M-1+J\), and define the individual
prime-power feature maps by
\[
 (\mathcal D_nq)(x)=
 \left(\frac{\Lambda(n)}{\sqrt n}K(x)K(x-\log n)\right)^{1/2}
 \bigl(q(x)-q(x-\log n)\bigr),
 \qquad n=p^k.
\]
Set
\[
 \mathcal D_{X,Y}=\bigoplus_{\substack{n=p^k\\X<n\le Y}}\mathcal D_n.
\]
This is a square-root feature map for the finite omitted prime-power
block:
\[
 \|\mathcal D_{X,Y}q\|^2
 =\sum_{\substack{n=p^k\\X<n\le Y}}
 \frac{\Lambda(n)}{\sqrt n}\mathcal J_{\log n}(q,q).
\tag{9.25ap}
\]
Realize its observation space as \(L^2(\Omega_Y,\nu_Y)\), where
\(\Omega_Y\) is the disjoint union of the channel copies of \(\mathbb R\).
The one-particle exterior space is
\[
 \mathcal H_Y=\mathbb C^{M-1}\oplus L^2(\Omega_Y,\nu_Y)
 =L^2(\widetilde\Omega_Y,\widetilde\nu_Y),
\]
where \(\widetilde\Omega_Y\) adjoins \(M-1\) counting-measure atoms.
We use the normalized exterior convention
\[
 \|\Phi\|_{[r]}^2=\frac1{r!}
 \int_{\widetilde\Omega_Y^r}|\Phi|^2\,\dd\widetilde\nu_Y^r.
\]
Abbreviate
\[
 A=\widehat A,\qquad
 U=\mathcal D_{X,Y}\Phi,\qquad
 W=\mathcal D_{X,Y}\Psi,\qquad
 v=\mathcal D_{X,Y}q_{M,J,X}^* .
\]
The radical basis and finite observability give \(W^*W\succ0\).  Set
\[
 \Pi_Y=I-W(W^*W)^{-1}W^*,\qquad
 z=\Pi_Yv,\qquad
 B=\Pi_YUA^{-1}U^*\Pi_Y .
\tag{9.25aq}
\]
The finite adaptive gain is then
\[
 \boxed{
 G_Y:=\mathfrak C_{M,J;X,Y}
 =\langle z,(I+B)^{-1}z\rangle
 =\frac{\tau_{d+1}(Y)}{\tau_d(Y)}.}
\tag{9.25ar}
\]

\begin{theorem}[Exact signed amplitude transfer]
\label{thm:signed-amplitude-transfer}
Consider a negative row with \(\widehat A\succ0\).
Let \(\psi\in\ker W^*\) and
\(\beta_\psi=\langle\psi,v\rangle\ne0\).  Define
\[
 g_\psi=(-A^{-1/2}U^*\psi,\psi),
 \qquad
 h_\psi=\frac{(-1)^d\sqrt{\delta^{\flat}_{M,J,X}}}
          {\overline{\beta_\psi}}\,g_\psi .
\tag{9.25as}
\]
For columns \(c_1,\dots,c_r\in\mathcal H_Y\) put
\(\Delta_r(\omega_1,\dots,\omega_r)=\det[c_j(\omega_i)]_{i,j\le r}\).
Let \(\Delta_d\) use the \(d\) nuisance columns
\[
 (A^{1/2}e_i,Ue_i)\ (i\le M-1),\qquad (0,We_a)\ (a\le J),
\]
and let \(\Delta_{d+1}\) append the column \((0,v)\), with the integrated
variable \(\omega\) in the first row.  By Andreief's identity
\cite{Andreief1886}, \(\|\Delta_r\|_{[r]}^2\) is the Gram determinant of
\(c_1,\dots,c_r\); thus \(\|\Delta_d\|_{[d]}^2=\tau_d(Y)\) and
\(\|\Delta_{d+1}\|_{[d+1]}^2=\tau_{d+1}(Y)\).  Then
\[
 \boxed{
 \int\overline{h_\psi(\omega)}
 \Delta_{d+1}(\omega,\boldsymbol\omega)\,
 \dd\widetilde\nu_Y(\omega)
 =\sqrt{\delta^{\flat}_{M,J,X}}\Delta_d(\boldsymbol\omega).}
\tag{9.25at}
\]
The associated exterior interior-product operator has norm
\[
 \boxed{
 \|T_{h_\psi}\|_{\rm op}^2
 =\delta^{\flat}_{M,J,X}
 \frac{\|\psi\|^2+\|A^{-1/2}U^*\psi\|^2}
  {|\langle\psi,v\rangle|^2}.}
\tag{9.25au}
\]
\end{theorem}

\begin{proof}
Expand \(\Delta_{d+1}\) along its first row and integrate against
\(\overline{g_\psi}\).  The pairing of \(g_\psi\) with an old-mode column
is \(-\langle U^*\psi,e_i\rangle+\langle U^*\psi,e_i\rangle=0\), with a
radical column it is \(\langle W^*\psi,e_a\rangle=0\), and with \((0,v)\)
it is \(\beta_\psi\).  Only the cofactor of the last column survives,
which is \((-1)^d\Delta_d\); multiplying by the factor in (9.25as) gives
(9.25at).  Contraction by a
one-particle vector on an exterior power has norm \(\|h_\psi\|\), which
gives (9.25au).
\end{proof}

By (9.25ai) and (9.25ar), \(G_Y>0\), so \(z\ne0\).  The choice
\(\psi=z\) gives the matched-filter certificate
\[
 Q_1=\frac{m_0^2}{m_0+m_1},
 \qquad m_j=\langle z,B^jz\rangle .
\tag{9.25av}
\]
It is sufficient, but it need not detect a crossing already present in
\(G_Y\).

For \(k\ge1\), define
\[
 \mathcal K_k=\operatorname{span}\{z,Bz,\ldots,B^{k-1}z\},
\]
\[
 H_k=[m_{i+j}+m_{i+j+1}]_{i,j=0}^{k-1},\qquad
 b_k=(m_0,\ldots,m_{k-1})^{\mathsf T},\qquad
 Q_k=b_k^*H_k^\dagger b_k .
\tag{9.25aw}
\]

\begin{theorem}[Adaptive Krylov hierarchy]
\label{thm:adaptive-krylov}
One has
\[
 \boxed{
 Q_k=\sup_{0\ne\psi\in\mathcal K_k}
 \frac{|\langle\psi,z\rangle|^2}
  {\langle\psi,(I+B)\psi\rangle},
 \qquad
 Q_1\le Q_2\le\cdots\le G_Y.}
\tag{9.25ax}
\]
If the spectral measure of \(B\) seen from \(z\) has \(r\) support
points, then \(Q_r=G_Y\).
\end{theorem}

\begin{proof}
Writing \(\psi=\sum_{j<k}c_jB^jz\) converts the quotient into
\(|c^*b_k|^2/(c^*H_kc)\).  The finite Riesz formula gives \(Q_k\), and
nestedness gives monotonicity.  If \(c\in\ker H_k\), its Krylov vector is
zero, hence \(c^*b_k=0\); thus \(b_k\in\operatorname{ran}H_k\) and the
pseudoinverse formula is well-defined.  On an \(r\)-point support,
\((1+t)^{-1}\) has a polynomial interpolant of degree at most \(r-1\),
so \((I+B)^{-1}z\in\mathcal K_r\).
\end{proof}

\subsection{Second-level spectral dispersion}
\label{sec:q2-dispersion}

Put
\[
 N_2=m_0m_2-m_1^2,\qquad
 D_2=(m_0+m_1)(m_2+m_3)-(m_1+m_2)^2 .
\tag{9.25ay}
\]
When \(D_2>0\), direct inversion of \(H_2\) gives
\[
 \boxed{
 Q_2=Q_1+\frac{N_2^2}{(m_0+m_1)D_2}.}
\tag{9.25az}
\]
If \(D_2=0\), the spectral measure seen from \(z\) has one support point
and \(Q_2=Q_1=G_Y\).

Let \(d\rho(t)=d\langle z,E_B(t)z\rangle\).  Then
\[
 \boxed{
 N_2=\frac12\iint(s-t)^2\,\dd\rho(s)\dd\rho(t),}
\]
\[
 \boxed{
 D_2=\frac12\iint
 (s-t)^2(1+s)(1+t)\,\dd\rho(s)\dd\rho(t).}
\tag{9.25ba}
\]
Equivalently,
\[
 N_2=\|z\wedge Bz\|^2,\qquad
 D_2=\|(I+B)^{1/2}z\wedge(I+B)^{1/2}Bz\|^2 .
\tag{9.25bb}
\]
Thus, when \(Q_1<\delta^{\flat}_{M,J,X}\),
\[
 \boxed{
 N_2^2>
 (\delta^{\flat}_{M,J,X}-Q_1)(m_0+m_1)D_2}
\tag{9.25bc}
\]
is exactly equivalent to \(Q_2>\delta^{\flat}_{M,J,X}\).

For any two orthonormal observation packets \(\phi_0,\phi_1\), put
\[
 a_i=\langle\phi_i,z\rangle,\qquad
 b_i=\langle\phi_i,Bz\rangle .
\]
The corresponding exterior coordinate gives
\[
 N_2\ge|a_0b_1-a_1b_0|^2.                       \tag{9.25bd}
\]
Since \(D_2\le(1+\|B\|)^2N_2\), the signed two-packet condition
\[
 \boxed{
 |a_0b_1-a_1b_0|^2>
 (\delta^{\flat}_{M,J,X}-Q_1)(m_0+m_1)(1+\|B\|)^2}
\tag{9.25be}
\]
implies the finite crossing.

The literal ordinary-prime form uses one common radical and old-mode
regression.  Put
\[
 \zeta^*=-(W^*W)^{-1}W^*v,\qquad
 Z=-(W^*W)^{-1} W^*U,
\]
\[
 \widetilde q=q_{M,J,X}^*+\Psi\zeta^*,\qquad
 \widetilde\Phi=\Phi+\Psi Z,
\]
and, in the individual prime-power channels,
\[
 x_n=\mathcal D_n\widetilde q,\qquad
 F_n=\mathcal D_n\widetilde\Phi A^{-1/2},\qquad
 \xi=\sum_{\substack{n=p^k\\X<n\le Y}}F_n^*x_n .
\tag{9.25bf}
\]
Then
\[
 \boxed{
 N_2=\frac12
 \sum_{\substack{n=p^k,\ \ell=r^a\\X<n,\ell\le Y}}\iint
 \left|x_n(s)(F_\ell\xi)(t)
 -(F_n\xi)(s)x_\ell(t)\right|^2\,\dd s\,\dd t.}
\tag{9.25bg}
\]
All von Mangoldt weights and theta phases therefore remain coupled until
the final modulus square.

\begin{corollary}[Exact hierarchy formulation of the frontier]
\label{cor:krylov-frontier}
For every fixed negative anti-shorted row with \(\widehat A\succ0\),
\[
 \boxed{
 \mathfrak g_{M,J,X}>\delta^{\flat}_{M,J,X}
 \Longleftrightarrow
 \exists\,Y<\infty,\ \exists\,k<\infty:
 Q_k(Y)>\delta^{\flat}_{M,J,X}.}
\tag{9.25bh}
\]
\end{corollary}

\begin{proof}
The forward implication follows from augmented determinant exhaustion,
followed by finite Krylov exactness.  Conversely,
\(Q_k(Y)\le G_Y\le\mathfrak g_{M,J,X}\).
\end{proof}

\begin{remark}[Level two is a sufficient gate]
The depth in Corollary~\ref{cor:krylov-frontier} may depend on the finite
row.  The dispersion inequality (9.25bc) is the first explicit
strengthening of the matched-filter certificate \(Q_1\), not a replacement for the exact surplus.  For
example,
\[
 B=\operatorname{diag}(1,2,3),\qquad
 z=\sqrt{12/13}(1,1,1),\qquad
 \delta=649/650
\]
gives \(G_Y=1>\delta\), but \(Q_2=648/650<\delta\).  Failure at level
two is therefore inconclusive.  If \(G_Y>\delta\), finite Krylov
exactness guarantees that some higher row-dependent level detects the
crossing.
\end{remark}

\begin{problem}[Physical strict surplus]
\label{prob:physical-strict-surplus}
For every negative anti-shorted row with \(\widehat A\succ0\) required in the cofinal exhaustion,
prove, with the literal theta translations and ordinary von Mangoldt
weights,
\[
 \boxed{\mathfrak g_{M,J,X}>\delta^{\flat}_{M,J,X}.}
\tag{9.25bi}
\]
By Theorem~\ref{thm:augmented-exhaustion}, this is equivalent to one
finite inequality (9.25ao), and by
Corollary~\ref{cor:krylov-frontier} to one finite Krylov certificate of
row-dependent depth.  By (9.25an) it is equivalent to
\(\mathcal A_\infty(q^*_{M,J,X},q^*_{M,J,X})>\mathfrak l_{M,J,X}\).
\end{problem}

\subsection{Heat-core exhaustion of the radical complement}
\label{subsec:heat-core}

The elementary zero modes provide the explicit orbit coordinate used in
the finite calculations above.  Their completeness in the weighted form
topology is a separate spectral-synthesis statement.  It is not needed for
the global passage.  We now construct a form core directly from the closed
ordinary-prime--Gamma form.

Let \(\mathcal R:=\overline{\bigcup_J\mathfrak R_J}^{L^2(\mu_K)}\) and
\[
 \mathscr C=(\mathbf1\oplus\mathcal R)^\perp,
 \qquad
 \mathcal Q_{\mathscr C}=D(\mathscr E_K)\cap\mathscr C.
\tag{9.25bj}
\]
Here \(\mathscr E_K\) is the closed form constructed as follows.  The
jump form is the squared norm of the difference operator
\[
\begin{aligned}
 \mathcal Dq={}&
 \left(
 \sqrt{g_\Gamma(u)K(x)K(x-u)}\{q(x)-q(x-u)\}
 \right)_{(u,x)\in(0,\infty)\times\mathbb R}
 \\
 &\oplus
 \left(
 \sqrt{\frac{\Lambda(n)}{\sqrt n}K(x)K(x-\log n)}
 \{q(x)-q(x-\log n)\}
 \right)_{n=p^a},
\end{aligned}
\tag{9.25bk}
\]
from \(L^2_{\rm even}(\mu_K)\) to the corresponding Hilbert direct sum.
If \(q_\ell\to q\) and \(\mathcal Dq_\ell\to Y\), one subsequence converges
almost everywhere with respect to Lebesgue measure because \(d\mu_K\) has
a strictly positive density.  Fubini gives simultaneous convergence at
\(x-u\), and an almost-everywhere subsequence in the disjoint-union output
space identifies \(Y=\mathcal Dq\).  Thus \(\mathcal D\) is closed.
Moreover \(C_{c,\rm even}^\infty(\mathbb R)\subset D(\mathcal D)\): near
\(u=0\), \(g_\Gamma(u)=O(u^{-1})\) and the squared difference is \(O(u^2)\),
while for large \(u\) the theta overlap is \(O_q(e^{-c_qe^{2u}})\).
At \(u=\log n\) this is summable against
\(\Lambda(n)/\sqrt n\).  The restriction of the closed operator
\(\mathcal D\) to \(C_{c,\rm even}^\infty(\mathbb R)\) is therefore
closable; let \(\mathscr E_K(q,s)=\langle\mathcal Dq,\mathcal Ds\rangle\)
on the domain \(D(\mathscr E_K)\) of its closure.  This is a densely
defined closed form; let \(L\ge0\) be its associated self-adjoint
operator.  The functions \(\mathbf1\), \(r_j\) and \(\chi_z\) (with
\(\lvert\operatorname{Im}z\rvert<\frac12\)) lie in \(D(\mathscr E_K)\):
for a cutoff \(\varphi_R(x)=\varphi(x/R)\) with \(\varphi\in
C_{c,\rm even}^\infty\) equal to \(1\) near \(0\), the functions
\(\varphi_Rq\) converge to \(q\) in \(L^2(\mu_K)\) and
\(\mathcal D(\varphi_Rq)\to\mathcal Dq\) by dominated convergence,
since \(q\) and \(q'\) grow at most exponentially while
\(K(x)\ll e^{5|x|}e^{-\pi e^{2|x|}}\).

Equation (9.15), extended by continuity from
\(f\in C_{c,\rm even}^\infty\) to \(f\in D(\mathscr E_K)\), and
\(\mathcal D\mathbf1=0\) give, by the representation theorem, the strong
domain identities (conditional on the ground-state identity, as is
(9.15))
\[
 L\mathbf1=0,
 \qquad
 r_j\in D(L),\qquad Lr_j=\frac12r_j,
\tag{9.25bl}
\]
hence also \(Lr=\frac12r\) for \(r\in\mathcal R\).  Therefore \(\mathbf1\oplus\mathcal R\) and \(\mathscr C\) reduce \(L\).
From here to the end of this section the ground-state identity stated
after (9.4) is assumed; it is used through (9.25bl) to make
\(\mathscr C\) a reducing subspace of \(L\).  Put
\[
 S_{\mathscr C}=L|_{\mathscr C}+\frac12I\ge\frac12I.
\tag{9.25bm}
\]

Choose a sequence \((g_j)_{j\ge1}\subset\mathscr C\) whose algebraic span
is Hilbert-dense and define
\[
 V_M^{\rm heat}
 =\operatorname{span}\{e^{-S_{\mathscr C}/k}g_j:1\le j,k\le M\}.
\tag{9.25bn}
\]

\begin{theorem}[Heat-core Galerkin exhaustion]
\label{thm:heat-core}
Assume the ground-state identity.  The spaces \(V_M^{\rm heat}\) are nested finite-dimensional subspaces of
\(\mathcal Q_{\mathscr C}\), every one of their vectors belongs to
\(\cap_{m\ge0}D(S_{\mathscr C}^m)\), and
\[
 \boxed{
 \overline{\bigcup_MV_M^{\rm heat}}^{\,
 (\frac12\|\cdot\|_{\mu_K}^2+\mathscr E_K)^{1/2}}
 =\mathcal Q_{\mathscr C}.}
\tag{9.25bo}
\]
For every fixed finite-dimensional
\(E\subset D(\mathscr E_K)\), the omitted Euler tail satisfies
\[
 \boxed{
 \left\|N_E^{-1/2}(\mathcal T_X)_EN_E^{-1/2}\right\|
 \longrightarrow0.}
\tag{9.25bp}
\]
Finally, if a finite omitted block contains two distinct primes \(p,r\),
its Gram form is positive definite on every finite centered space
\(E\oplus\mathfrak R_J\) with \(E\subset\mathscr C\).
\end{theorem}

\begin{proof}
Spectral calculus gives, for \(t>0\),
\[
 \|S_{\mathscr C}^{m}e^{-tS_{\mathscr C}}\|
 \le\left(\frac{m}{et}\right)^m,
 \qquad
 \|S_{\mathscr C}^{1/2}e^{-tS_{\mathscr C}}\|
 \le(2et)^{-1/2}.
\tag{9.25bq}
\]
For \(q\in D(S_{\mathscr C}^{1/2})\), dominated convergence in its spectral
measure gives
\[
 \|S_{\mathscr C}^{1/2}(e^{-tS_{\mathscr C}}-I)q\|\longrightarrow0.
\tag{9.25br}
\]
Given \(\varepsilon>0\), first choose \(k\) so that (9.25br), with
\(t=1/k\), is smaller than \(\varepsilon/2\).  Approximate \(q\) in the
Hilbert norm by a finite linear combination \(g\) of the \(g_j\), with
\(\|q-g\|<(\varepsilon/2)\sqrt{2et}\).  Equations (9.25bq)--(9.25br)
then give
\[
 \|S_{\mathscr C}^{1/2}(q-e^{-tS_{\mathscr C}}g)\|<\varepsilon,
\]
which proves (9.25bo).

For (9.25bp), choose a basis \(e_1,\ldots,e_d\) of \(E\).  Positivity and
monotone convergence give
\(\mathcal T_X(e_i,e_i)\downarrow0\), while Cauchy--Schwarz for the tail
form gives
\[
 |\mathcal T_X(e_i,e_j)|^2
 \le\mathcal T_X(e_i,e_i)\mathcal T_X(e_j,e_j).
\]
Thus every entry, and hence the normalized finite matrix, tends to zero.

For the last assertion, vanishing of the two prime atoms would give the
two periods \(\log p\) and \(\log r\).  Their quotient is irrational,
since a rational relation would imply \(p^a=r^b\).  Strong continuity of
translations in \(L^2_{\rm loc}\) then makes every real number a period,
so the vector is constant almost everywhere.  Centeredness makes it zero.
The finite Gram matrix is therefore positive definite.
\end{proof}

The hybrid spaces
\[
 V_M^{\rm hyb}=V_M^{\rm heat}+V_M
\tag{9.25bs}
\]
are also a form core.  They retain every explicit zero-orbit diagnostic
used above, without asserting that the canonical zero modes alone are
dense.  One-channel observability and the local scalar Christoffel bounds
remain statements about the canonical summand.  In the global closure,
finite tail blocks must contain two distinct primes, and the physical
surplus must be proved on the heat or hybrid rows.

\subsection{The compensated frontier and closure}

The strict frontier (9.12) excludes a completed threshold null vector and
can therefore be stronger than mere nonnegativity.  The exact
vanishing-tolerance formulation is the following.  Choose
\[
 M\longmapsto(J(M),X(M),\varepsilon_M),
\quad J(M),X(M)\to\infty,\quad\varepsilon_M\downarrow0,
\tag{9.26}
\]
and require
\[
 \boxed{
 n_-\!\left[
 \mathbf A_{M,J(M),X(M)}
 +\varepsilon_M
 \begin{pmatrix}\mathbf N_M&0\\0&0\end{pmatrix}
 \right]=J(M).}
\tag{9.27}
\]
Equivalently,
\[
 \boxed{
 \mathbf T^{VV}
 -\mathbf T^{VR}(\mathbf T^{RR})^{-1}\mathbf T^{RV}
 \preceq\mathbf H_M+\varepsilon_M\mathbf N_M.}
\tag{9.28}
\]
On the canonical zero-mode summand the cutoff part of (9.28) has the
explicit bound \(C_Me^{-cX(M)}\mathbf N_M\).  On every heat or hybrid row,
Theorem~\ref{thm:heat-core} gives the qualitative bound: after the finite
row is fixed, \(X(M)\) may be enlarged until the normalized tail is at most
\(\varepsilon_M\).  What remains is the inequality (9.28) itself.

\begin{theorem}[Cofinal inertia closure]
\label{thm:cofinal-inertia-closure}
Assume the ground-state identity stated after (9.4).  Let \(V_M\) be (9.25bn) or (9.25bs), with \(\mathbf N_M\),
\(\mathbf H_M\) and \(\mathbf A_{M,J,X}\) defined by (9.6) and (9.17)
for this \(V_M\), and assume (9.27) for some \((J(M),X(M),\varepsilon_M)\)
with \(\varepsilon_M\to0\).  Then
\[
 \mathcal A_\infty(q,q)\ge0
\tag{9.29}
\]
for every \(q\in\mathcal Q_{\mathscr C}\), hence on
\(D(\mathscr E_K)\cap\mathcal H_0\); by (7.5) this is Weil positivity on
even tests, which gives RH together with the corresponding inequality for
odd tests (Theorem~\ref{thm:weil}).
\end{theorem}

\begin{proof}
By (9.19), with \(\mathbf H_M\) replaced by
\(\mathbf H_M+\varepsilon_M\mathbf N_M\), (9.27) is equivalent to
(9.28).  The tail Schur complement
\(\mathbf T^{VV}-\mathbf T^{VR}(\mathbf T^{RR})^{-1}\mathbf T^{RV}\) is
positive semidefinite, being a Schur complement of the positive
semidefinite tail Gram matrix; hence (9.28) gives
\(\mathbf H_M\succeq-\varepsilon_M\mathbf N_M\).

Now fix \(q\in V_m\).  For every \(M\ge m\), \(q\in V_M\), and the
previous inequality gives
\[
 \mathcal A_\infty(q,q)
 \ge-\varepsilon_M\|q\|_{\mu_K}^2.
\]
Letting \(M\to\infty\) gives \(\mathcal A_\infty(q,q)\ge0\) on
\(\cup_MV_M\).  Since \(\mathcal A_\infty\) is continuous in the closed
form norm, (9.25bo) (for the hybrid spaces, the same density together
with \(V_M^{\rm heat}\subset V_M^{\rm hyb}\subset\mathcal Q_{\mathscr C}\))
extends the inequality to \(\mathcal Q_{\mathscr C}\).

By (9.15) and (9.25bl),
\(\mathcal A_\infty\) vanishes on \(\mathcal R\) and
\(\mathcal A_\infty(r,f)=0\) for \(r\in\mathcal R\) and \(f\) in the form
domain, so \(\mathcal A_\infty\ge0\) on
\(D(\mathscr E_K)\cap\mathcal H_0=\mathcal Q_{\mathscr C}\oplus
\mathcal R\).  For even \(\rho\in C_c^\infty(\mathbb R)\), put
\(\rho_0=\rho-\int\rho\,\dd\mu_K\in D(\mathscr E_K)\cap\mathcal H_0\); since
\(\mathcal D\) annihilates constants, (7.5) and the ground-state identity
give \(\QW(K\rho,K\rho)=\mathcal A_\infty(\rho_0,\rho_0)\ge0\).  This is even Weil
positivity; as in the proof of the gate proposition in Section~7, it gives
RH by Theorem~\ref{thm:weil} together with the inequality for odd tests.
\end{proof}

This section proves the identities (9.19), (9.25g), (9.25aa),
(9.25ai), (9.25ax) and Theorem~\ref{thm:adaptive-krylov}; (9.15),
(9.25bl), Theorem~\ref{thm:heat-core} and
Theorem~\ref{thm:cofinal-inertia-closure} are conditional on the
ground-state identity.  The inequality (9.25bi) is not proved; by
(9.25an) it is equivalent to
\(\mathcal A_\infty(q^*_{M,J,X},q^*_{M,J,X})>\mathfrak l_{M,J,X}\), a
positivity statement for the completed form itself.  If (9.25bi) holds
on every negative row (with \(X\) large enough that
\(\widehat A\succ0\)), then, since \(\mathfrak g_{M,J,X}>0\) by
(9.25ad) and the last assertion of Theorem~\ref{thm:heat-core} (the
infimum in (9.25ad) is attained at a nonzero vector, on which the omitted
tail is strictly positive), the completed last pivot (9.25aa) is positive on every row, so
by induction on \(M\) and (9.25g) \(\mathbf H_M\succ0\) on each \(V_M\);
by (9.25bp) and \(\mathbf T^{RR}\succ0\) there is \(X(M)\) for which
(9.28) holds with \(\varepsilon_M=0\), hence (9.27).

\begin{remark}[Strict and compensated targets]
The strict statement requested by the finite-energy-level picture is
\[
 \boxed{X_*(M,J(M))<\infty\quad\text{for every }M.}
\tag{9.30}
\]
By the argument of Theorem~\ref{thm:exact-finite-repair}, applied to the
shorted tail
\(\mathbf T^{VV}-\mathbf T^{VR}(\mathbf T^{RR})^{-1}\mathbf T^{RV}\)
(strictly positive on \(V_M\) and tending to \(0\)), it is equivalent to
\(\mathbf H_M\succ0\).  The compensated statement (9.27) is
the threshold-stable version used in Theorem~\ref{thm:cofinal-inertia-closure}.  A proof of (9.30) on a
form-core exhaustion proves (9.27) immediately; a proof of (9.27) need
not exclude additional threshold null states at finite \(M\).
\end{remark}

\section{Tate boundary geometry and global polarization}
\label{sec:polarization-geometry}

The finite-inertia analysis above isolates the missing sign but does not
explain where that sign could come from.  This section records a second,
source-level construction.  It starts from finite root graphs at the
ordinary primes, constructs the Gamma and polar determinant at infinity,
and joins them before completion.  The construction produces the correct
local masses, parity, boundary cancellation and scale covariance.  It also
determines the spectral type of the ordinary Hilbert completion.

The distinction maintained throughout is
\[
 \boxed{
 \text{positive local/source polarization}
 \ \not\Longrightarrow\ 
 \text{positive faithful polarization of }H^1_{\rm CCM}.}
 \tag{10.1}
\]
The left side will be constructed below.  The most natural absolutely
continuous realization of the arrow in (10.1) will be disproved.

\subsection{Finite root geometry and the literal von Mangoldt mass}

For \(N\ge2\), let \(R_N=\mathbb Z/N\mathbb Z\), let \(T_N\) be
translation by one on \(\ell^2(R_N)\), and put
\[
 \Delta_N=2I-T_N-T_N^*.
 \tag{10.2}
\]
Its form is the positive graph energy
\[
 \langle f,\Delta_Nf\rangle
 =\sum_{j\in R_N}|f(j+1)-f(j)|^2.
 \tag{10.3}
\]
The kernel consists of the constants.  We write \(\det'\Delta_N\) for
the product of the nonzero eigenvalues.

\begin{theorem}[Root torsion--incidence identity]
\label{thm:root-torsion}
For every \(N\ge2\),
\[
 \det'\Delta_N=N^2.
 \tag{10.4}
\]
Fix \(m\ge2\), a prime \(p\), and \(k\ge1\).  Then
\[
 \frac12\log
 \frac{\det'\Delta_{mp^k}}{\det'\Delta_{mp^{k-1}}}
 =\log p.
 \tag{10.5}
\]
If \(H_{m,k}=p^kR_m\subset R_{mp^k}\), and
\[
 \omega_{mp^k}=(mp^k)^{-1/2}\mathbf1_{R_{mp^k}},
 \qquad
 \eta_{m,k}=m^{-1/2}\mathbf1_{H_{m,k}},
\]
then
\[
 \langle\omega_{mp^k},\eta_{m,k}\rangle=p^{-k/2}.
 \tag{10.6}
\]
Consequently
\[
 \boxed{
 \left(\frac12\log
 \frac{\det'\Delta_{mp^k}}{\det'\Delta_{mp^{k-1}}}\right)
 \langle\omega_{mp^k},\eta_{m,k}\rangle
 =\frac{\Lambda(p^k)}{\sqrt{p^k}}.}
 \tag{10.7}
\]
\end{theorem}

\begin{proof}
The Fourier characters of \(R_N\) diagonalize \(\Delta_N\), with
eigenvalues \(|1-e^{2\pi ij/N}|^2\).  Hence
\[
 \det'\Delta_N
 =\left|\prod_{j=1}^{N-1}(1-e^{2\pi ij/N})\right|^2=N^2,
\]
which gives (10.4) and (10.5).  The subgroup \(H_{m,k}\) has \(m\)
elements, so the normalized overlap is
\(m/(\sqrt{mp^k}\sqrt m)=p^{-k/2}\).  Multiplying by (10.5)
proves (10.7).
\end{proof}

The two factors in (10.7) have different meanings.  The logarithm is a
relative determinant increment along the root tower, whereas
\(p^{-k/2}\) is a normalized incidence number.  The distinction prevents
the incorrect replacement of the weight of the layer \(p^k\) by
\(k\log p\).

The finite root levels are polarized before any global sign is considered.
On \(H_N^0=\mathbf1^\perp\subset\ell^2(R_N)\), set
\[
 \begin{aligned}
 J_N(x,y)&=(-y,x),\\
 g_N((x,y),(u,v))
 &=\langle\Delta_Nx,u\rangle+\langle\Delta_Ny,v\rangle,\\
 \Omega_N(z,w)&=-g_N(z,J_Nw).
 \end{aligned}
 \tag{10.8}
\]
Then \(g_N\) is positive definite, \(J_N^2=-I\), and \(\Omega_N\) is
alternating and nondegenerate.  If \(M\mid N\), normalized pullback along
\(R_N\to R_M\) is a polarized isometry and the determinant on the new-root
complement is \((N/M)^2\).  Thus (10.8) is functorial under the root
coverings which produce (10.5).

\subsection{Poisson holonomy, heat lift and local unitary triviality}

The coefficient \(p^{-k/2}\) also has an exact positive realization as a
unitary moment.  For \(0<r<1\), let
\[
 \dd\mu_r(\theta)
 =\frac{1-r^2}{1-2r\cos\theta+r^2}\frac{\dd\theta}{2\pi}.
\tag{10.8a}
\]
The Fourier expansion of the Poisson kernel gives
\[
 \int_{-\pi}^{\pi}e^{ik\theta}\dd\mu_r(\theta)=r^{|k|}.
\tag{10.8b}
\]
Taking \(r=p^{-1/2}\) realizes every winding of the prime tower without
analytic continuation.

Let \(A_{p,\theta}\) be the nonnegative Laplacian on the circle of length
\(\ell_p=\log p\), with boundary condition
\(f(x+\ell_p)=e^{i\theta}f(x)\).  The method-of-images kernel is
\[
 K_{p,\theta,t}(x,y)
 =\frac1{\sqrt{4\pi t}}
  \sum_{k\in\mathbb Z}
  e^{-(x-y+k\ell_p)^2/(4t)}e^{-ik\theta}.
\tag{10.8c}
\]
It satisfies the exact heat composition law before any winding sector is
removed.  Let \(A_p=\int^\oplus A_{p,\theta}\,\dd\mu_{p^{-1/2}}(\theta)\) and
\(\tau_p(T)=\int\Tr(T_\theta)\,\dd\mu_{p^{-1/2}}(\theta)\) for decomposable
\(T=\int^\oplus T_\theta\).  Integrating the trace of (10.8c) against
\(\dd\mu_{p^{-1/2}}\) in (10.8a) yields
\[
 \boxed{
 \tau_p(e^{-tA_p})
 =\frac{\ell_p}{\sqrt{4\pi t}}
  \left(1+2\sum_{k\ge1}p^{-k/2}
  e^{-(k\ell_p)^2/(4t)}\right).}
\tag{10.8d}
\]
The nonzero-winding relative trace is therefore exactly
\[
 \tau_p(e^{-tA_p})-\frac{\log p}{\sqrt{4\pi t}}
 =\frac{\log p}{\sqrt{\pi t}}
  \sum_{k\ge1}p^{-k/2}e^{-(k\log p)^2/(4t)}.
\tag{10.8e}
\]
This is the complete Gaussian prime tower in the heat explicit formula.

The local realization has a rigidity which is important globally.  The
Poisson density is the modulus square of the outer function
\[
 a_r(e^{i\theta})=\frac{\sqrt{1-r^2}}{1-re^{i\theta}}.
\tag{10.8f}
\]
Multiplication by \(a_r\) is a unitary map
\(L^2(\mathbb T,\mu_r)\to L^2(\mathbb T,\dd\theta/2\pi)\) and commutes
with every decomposable fiber operator.  Hence the weighted and unweighted
local Laplacians are unitarily equivalent.  The prime changes the normal
weight used to read the heat diagonal; it does not create a new local
spectrum.  Moreover the winding-zero sector cannot be deleted while
preserving composition, because two nonzero windings can add to zero.
Thus the finite quantity in (10.8e) is a relative semifinite trace, not the
ordinary trace of a reduced semigroup.

There is an equivalent precision-space formulation.  On
\(\ell^2(\mathbb N_0)\), let
\[
 (K_px)_{h}=\sum_{k\ge0}p^{-|h-k|/2}x_k.
\]
With \(a=p^{-1/2}\), define
\[
 (B_ax)_0=x_0,
 \qquad
 (B_ax)_{h+1}=\frac{x_{h+1}-ax_h}{\sqrt{1-a^2}}.
\tag{10.8g}
\]
Direct triangular inversion gives
\[
 \boxed{
 K_p^{-1}=B_a^*B_a,
 \qquad
 \langle K_p^{-1}x,x\rangle
 =|x_0|^2+\frac1{1-p^{-1}}
  \sum_{h\ge0}|x_{h+1}-p^{-1/2}x_h|^2.}
\tag{10.8h}
\]
Thus the Poisson coefficient is also the Green kernel of a positive local
first-order arithmetic connection.  The local positivity is exact.  Heuristically, since
\(\sum_p p^{-1}=\infty\), one does not expect the infinite product of these
changes of coordinates to be a bounded perturbation of a reference product;
this is not proved here.
The matched Gamma--Euler boundary must therefore be formed before the
global norm or trace is taken.

\subsection{The Gamma spin and the polar determinant}

On \(\mathcal H_\Gamma=\ell^2(\mathbb N_0)\), define the positive operator
\[
 N_\Gamma e_m=\left(2m+\frac12\right)e_m.
\tag{10.9}
\]
Its heat trace is
\[
 \Tr e^{-uN_\Gamma}=\frac{e^{-u/2}}{1-e^{-2u}}.
\tag{10.10}
\]
Zeta regularization of \(N_\Gamma+s-1/2\) gives the archimedean
determinant directly from this source operator.

\begin{theorem}[Gamma--polar determinant]
\label{thm:gamma-polar-det}
For \(\Re s>0\), with meromorphic continuation thereafter,
\[
 \det_\zeta(N_\Gamma+s-1/2)
 =\frac{\sqrt{2\pi}\,2^{1/2-s/2}}{\Gamma(s/2)}.
\tag{10.11}
\]
Let \(N_{\rm triv}\) have eigenvalues \(-1/2\) and \(1/2\).  Then
\[
 \boxed{
 \frac12s(s-1)\pi^{-s/2}\Gamma(s/2)
 =\sqrt\pi(2\pi)^{-s/2}
 \frac{\det((s-1/2)I-N_{\rm triv})}
      {\det_\zeta(N_\Gamma+s-1/2)}.}
\tag{10.12}
\]
\end{theorem}

\begin{proof}
The spectral zeta function is
\[
 \Tr(N_\Gamma+s-1/2)^{-w}
 =2^{-w}\zeta_H(w,s/2).
\]
Use
\(\zeta_H(0,a)=1/2-a\) and
\(\zeta_H'(0,a)=\log\Gamma(a)-\frac12\log(2\pi)\) to obtain
(10.11).  Since
\(\det((s-1/2)I-N_{\rm triv})=s(s-1)\), substitution gives
(10.12).
\end{proof}

The positive degree-one Gamma page is the real double of the form domain
of \(N_\Gamma^{1/2}\), with metric and complex structure
\[
 \begin{aligned}
 g_\Gamma((x,y),(u,v))
 &=\langle N_\Gamma^{1/2}x,N_\Gamma^{1/2}u\rangle
  +\langle N_\Gamma^{1/2}y,N_\Gamma^{1/2}v\rangle,\\
 J_\Gamma(x,y)&=(-y,x).
 \end{aligned}
\tag{10.13}
\]
The polar plane remains in degrees zero and two.  It is not inserted as a
positive degree-one summand.  This keeps the parity \(H^0-H^1+H^2\) that, in the
function-field analogy, accounts for the sign of the prime terms in the
explicit formula; this heuristic is central to the geometric interpretation in
\S\ref{subsec:open-programs}.

\subsection{Matched boundary cancellation and the full finite part}

Let \(\mathscr K\) be a fixed separable Hilbert space (in (10.20) it is realized as \(\mathscr K_C\)).  For each
prime power set
\[
 w_{p,k}=\frac{\log p}{p^k}.
 \tag{10.14}
\]
The restricted-product boundary is decomposed into a common component
\(F\in\mathscr K\) and local deviations \(r=(r_{p,k})\), where
\[
 \sum_{p,k}w_{p,k}\|r_{p,k}\|
 +\left(\sum_{p,k}w_{p,k}\|r_{p,k}\|^2\right)^{1/2}<\infty.
 \tag{10.15}
\]
Thus the local row is \(F_{p,k}=F+r_{p,k}\).  For \(s>1/2\), put
\[
 C_s=\sum_p(\log p)p^{-2s},
 \qquad
 \kappa_\infty=\gamma_E+\sum_p\sum_{k\ge2}w_{p,k}.
 \tag{10.16}
\]
The Laurent expansion \(-\zeta'(\sigma)/\zeta(\sigma)=(\sigma-1)^{-1}-\gamma_E+O(\sigma-1)\)
(\(\gamma_E\) is Euler's constant)
at \(\sigma=2s\to1\), after subtraction of the terms \(k\ge2\), which
converge absolutely and uniformly near \(\sigma=1\), gives
\(C_s=(2s-1)^{-1}-\kappa_\infty+o(1)\) as \(s\downarrow1/2\).  Writing
\(\operatorname{FP}_{s\downarrow1/2}\) for the constant term in this
expansion (in the variable \(2s-1\)), we obtain
(numerically \(\kappa_\infty\approx1.3326\))
\[
 \operatorname{FP}_{s\downarrow1/2}C_s=-\kappa_\infty.
 \tag{10.17}
\]

For \(x=(F,r)\), \(y=(G,z)\), consider the jointly regularized form
\[
 \begin{aligned}
 \mathfrak Q_s(x,y)
 &=\sum_p(\log p)p^{-2s}
   \langle F+r_{p,1},G+z_{p,1}\rangle\\
 &\quad+\gamma_E\langle F,G\rangle
 +\sum_p\sum_{k\ge2}w_{p,k}
   \langle F+r_{p,k},G+z_{p,k}\rangle.
 \end{aligned}
 \tag{10.18}
\]

\begin{theorem}[Exact matched finite part]
\label{thm:matched-finite-part}
The finite part of (10.18) exists and equals
\[
 \boxed{
 \operatorname{FP}_{s\downarrow1/2}\mathfrak Q_s(x,y)
 =\sum_{p,k}w_{p,k}
 \bigl(
 \langle r_{p,k},G\rangle
 +\langle F,z_{p,k}\rangle
 +\langle r_{p,k},z_{p,k}\rangle
 \bigr).}
 \tag{10.19}
\]
The series is absolutely convergent and defines a continuous form on the
space in (10.15).  In particular the common generic diagonal has zero
finite part.
\end{theorem}

\begin{proof}
After expanding (10.18), the coefficient of \(\langle F,G\rangle\) is
\(C_s+\kappa_\infty\), whose finite part is zero by (10.17).  Every
residual coefficient converges to \(w_{p,k}\).  The first summand in
(10.15) controls the common--residual terms and weighted
Cauchy--Schwarz controls the residual--residual term, so dominated
convergence gives (10.19).
\end{proof}

Theorem~\ref{thm:matched-finite-part} is stronger than cancellation of
two scalar divergences.  It retains every independent local deviation with
its literal von Mangoldt weight.  Hence Gamma, the pole and the prime
layers are joined before taking an absolute square or a quotient.

\subsection{Covariant mixing of all prime-power charges}

The local row of charge \(\ell_{p,k}=k\log p\) has normalized scale
generator \(A-\ell_{p,k}\).  Adding rows with different charges without
transporting them is not equivariant.  The required transport is explicit.
In the spectral realization
\[
 \mathscr K_C=L^2(\mathbb R,w_C(\gamma)\dd\gamma),
 \qquad
 w_C(\gamma)=\frac1{2\pi(\gamma^2+1/4)},
 \qquad AF(\gamma)=\gamma F(\gamma),
 \tag{10.20}
\]
define
\[
 (S_\ell F)(\gamma)
 =\left(\frac{w_C(\gamma+\ell)}{w_C(\gamma)}\right)^{1/2}
  F(\gamma+\ell).
 \tag{10.21}
\]
Then \((S_\ell)_{\ell\in\mathbb R}\) is a strongly continuous unitary
group and
\[
 S_\ell(A-\ell I)=AS_\ell,
 \qquad
 S_\ell^*f(A)S_\ell=f(A-\ell I).
 \tag{10.22}
\]

For \(I_X=\{(p,k):p^k\le X\}\), let
\(\mathscr S_X=\bigoplus_{(p,k)\in I_X}(\mathscr K_C\oplus\mathscr K_C)\),
with vectors \(v=(x_{p,k},y_{p,k})_{(p,k)\in I_X}\).  Put
\[
 c_p=\frac{2\pi}{\log p},
 \qquad
 \alpha_{p,k}=\sqrt{\frac{\log p}{c_p}}p^{-k/2}.
 \tag{10.23}
\]
Equip \(\mathscr S_X\) with the metric
\(g_X(v,w)=\sum_{(p,k)\in I_X}\bigl(c_p\langle x_{p,k},x'_{p,k}\rangle
+c_p^{-1}\langle y_{p,k},y'_{p,k}\rangle\bigr)\) for \(w=(x',y')\), the
complex structure \(J_X(x,y)_{p,k}=(-c_p^{-1}y_{p,k},\,c_px_{p,k})\) (so
\(J_X^2=-I\) and \(J_X\) is \(g_X\)-orthogonal), and the scale action
\(U^X_t\) acting on the summand \((p,k)\) by
\(e^{it(A-\ell_{p,k})}\) in both components.
The covariant Tate boundary is
\[
 \begin{aligned}
 R_X^{\rm cov}v
 &=\sum_{p,k}\alpha_{p,k}S_{\ell_{p,k}}y_{p,k},\\
 R_X^{\rm cov}J_Xv
 &=\sum_{p,k}c_p\alpha_{p,k}S_{\ell_{p,k}}x_{p,k}.
 \end{aligned}
 \tag{10.24}
\]
It intertwines every charged source action with the one action generated
by \(A\): \(R_X^{\rm cov}U^X_t=e^{itA}R_X^{\rm cov}\) by (10.22).  Since
each \(S_\ell\) is unitary, with adjoints taken for \(g_X\),
\[
 R_X^{\rm cov}(R_X^{\rm cov})^*
 =R_X^{\rm cov}J_X(R_X^{\rm cov}J_X)^*
 =C_XI,
 \qquad
 C_X=\sum_{p^k\le X}\frac{\log p}{p^k}.
 \tag{10.25}
\]

For definiteness let the Gamma compliance symbol be
\[
 m_\Gamma(\gamma)=\Re\psi\Bigl(\frac14+\frac{i\gamma}2\Bigr)-\psi\Bigl(\frac14\Bigr)
 =\sum_{n\ge0}\Bigl(\frac1{n+1/4}-\frac{n+1/4}{(n+1/4)^2+\gamma^2/4}\Bigr),
\]
where \(\psi=\Gamma'/\Gamma\).  Each summand is nonnegative and vanishes
only at \(\gamma=0\), so \(m_\Gamma\) is real-analytic, nonnegative, and
its zero set is \(\{0\}\).  Let the complete Gamma compliance be
\[
 K_\Gamma=\kappa_\infty I+m_\Gamma(A)\succeq\kappa_\infty I.
 \tag{10.26}
\]
Shorting only the common archimedean boundary gives the positive metric
\[
 \begin{aligned}
 g_X^{\rm cov}(v,w)
 &=g_X(v,w)\\
 &\quad+\langle K_\Gamma^{-1}R_X^{\rm cov}J_Xv,
                         R_X^{\rm cov}J_Xw\rangle\\
 &\quad+\langle K_\Gamma^{-1}R_X^{\rm cov}v,
                         R_X^{\rm cov}w\rangle.
 \end{aligned}
 \tag{10.27}
\]
The diagonal row is
\[
 S_{\ell_{p,k}}^*K_\Gamma^{-1}S_{\ell_{p,k}}
 =\bigl(\kappa_\infty+m_\Gamma(A-\ell_{p,k})\bigr)^{-1},
 \tag{10.28}
\]
while two different rows retain the cross compliance
\[
 S_{\ell_i}^*K_\Gamma^{-1}S_{\ell_j}.
 \tag{10.29}
\]
Thus all prime rows use one common regression and all off-diagonal phases
survive.  Since each \(S_{\ell_{p,k}}\) is unitary, the diagonal terms
\(\langle F+r_{p,k},G+z_{p,k}\rangle\) of (10.18) are unchanged when both
arguments of the row \((p,k)\) are transported by \(S_{\ell_{p,k}}\); hence
(10.19) holds verbatim for the transported rows.  Cross terms between rows
of different charges are not of this form and do change.

Because \(K_\Gamma^{-1}\preceq\kappa_\infty^{-1}I\) and \(R_X^{\rm cov}\) is
bounded, \(g_X^{\rm cov}\) is equivalent to \(g_X\), and \(U^X_t\) is
\(g_X^{\rm cov}\)-unitary since \(K_\Gamma\) commutes with \(e^{itA}\).  For
\(X\le X'\) the inclusion \(\mathscr S_X\subset\mathscr S_{X'}\) is isometric
for \(g^{\rm cov}\), because \(R_{X'}^{\rm cov}\) restricts to
\(R_X^{\rm cov}\).  We write \(\mathscr P_\infty^{\rm cov}\) for the Hilbert
completion of \(\bigcup_X\mathscr S_X\) under \(g^{\rm cov}\), with the
unitary group \(U^{\rm cov}_t\) extending the \(U^X_t\).

\subsection{The global polarization target}

Let \(H^1_{\rm CCM}\) denote the separated resonant degree one suggested by
the CCM cyclic construction~\cite{CCM2025}, with weight-one scale action
\(\vartheta_t\) and descended alternating form \(\Omega\).  The properties
of this space that are used below are not proved here; we isolate them.

\medskip\noindent\textbf{Hypothesis (R).} \(H^1_{\rm CCM}\) is a
complex vector space with a linear action \(\vartheta_t\), \(t\in\R\), such
that for every nontrivial zero \(\rho\) of \(\zeta\) there is
\(u_\rho\neq0\) with \(\vartheta_tu_\rho=e^{t\rho}u_\rho\) for all \(t\).
\emph{Hypothesis (R) for critical-line zeros} is the same statement
restricted to zeros with \(\Re\rho=1/2\).  The eigenclasses \(u_\rho\)
are assumed to lie in every scale-invariant subspace used below.

\medskip\noindent The missing
global object can now be stated independently of the finite matrix model.

\begin{definition}[Global CCM polarization]
\label{def:global-polarization}
A global polarization is a real-linear operator \(J\) on a complex
\(\vartheta_t\)-invariant subspace \(\mathcal C\subset H^1_{\rm CCM}\)
containing all classes \(u_\rho\) of Hypothesis (R) such that
\[
 J^2=-I,
 \qquad
 J\vartheta_t=\vartheta_tJ,
 \qquad
 g(u,v):=\Omega(u,Jv)
 \tag{10.30}
\]
is a positive definite real inner product on \(\mathcal C\), where
\(\Omega\) is a real-bilinear alternating form on \(\mathcal C\).  The normalized action
\(U_t=e^{-t/2}\vartheta_t\) must extend unitarily to its completion.
\end{definition}

This is the infinite-dimensional analogue of the polarized degree-one
page in the proof of the Riemann hypothesis for curves over finite fields,
but the required positivity there is supplied by correspondences and the
Hodge index theorem on a square, rather than by a formal Hilbert
completion.  Deninger's dynamical-cohomology program gives the same
weight-one template~\cite{Deninger1998}.

\begin{proposition}[Exact force of the polarization]
\label{prop:polarization-rh}
Under Hypothesis (R), a global CCM polarization in the sense of
Definition~\ref{def:global-polarization} implies RH.
\end{proposition}

\begin{proof}
Let \(\rho\) be a nontrivial zero.  By Hypothesis (R) there is a class
\(u_\rho\ne0\) with
\(\vartheta_tu_\rho=e^{t\rho}u_\rho\); then \(u_\rho\in\mathcal C\) and
\(g(u_\rho,u_\rho)>0\), so \(u_\rho\) is a nonzero vector of the completion.  Then
\[
 U_tu_\rho=e^{t(\rho-1/2)}u_\rho.
\]
Write \(\rho-1/2=\sigma+i\tau\).  Then
\(U_tu_\rho=e^{t\sigma}(\cos(t\tau)u_\rho+\sin(t\tau)\,iu_\rho)\).  The
vectors \(u_\rho\) and \(iu_\rho\) are \(\mathbb R\)-linearly independent,
so on the real plane they span the norm \(\|\cdot\|_g\) is bounded above and
below on the compact curve \(\cos\phi\,u_\rho+\sin\phi\,iu_\rho\) by
constants \(0<c_1\le c_2\).  Since \(g(U_tu_\rho,U_tu_\rho)=g(u_\rho,u_\rho)\)
for all real \(t\), we get
\(c_1e^{t\sigma}\le\|u_\rho\|_g\le c_2e^{t\sigma}\) for all \(t\), hence
\(\sigma=0\), i.e.\ \(\Re\rho=1/2\).  Since
\(\rho\) was an arbitrary nontrivial zero, RH follows.
\end{proof}

Proposition~\ref{prop:polarization-rh} explains why positivity of the local
pages cannot be promoted by a formal completion: the promotion itself
contains the full missing sign.

\subsection{The absolutely continuous completion no-go}

The covariant system (10.20)--(10.29) has an ordinary Hilbert direct-limit
completion.  It is natural, positive and source-defined, but it has the
wrong spectral type.  The following operator-valued statement shows that the
Gamma compliance has dense range, so it creates no cokernel; the spectral
obstruction itself is Theorem~\ref{thm:ac-no-go}.

\begin{lemma}[Operator-valued dense-range criterion]
\label{lem:operator-dense-range}
Let \((X,\mu)\) be sigma finite, let \(\mathfrak h\) be separable, and let
\(M(x)\) be a measurable field of closed densely defined operators.  If
\(\mathcal M\) is the associated maximal decomposable operator on
\(L^2(X,\mu;\mathfrak h)\), then
\[
 \boxed{
 \overline{\operatorname{Ran}\mathcal M}
 =L^2(X,\mu;\mathfrak h)
 \quad\Longleftrightarrow\quad
 \ker M(x)^*=\{0\}\ \text{for almost every }x.}
 \tag{10.31}
\]
If \(M(x)^*M(x)\succeq\kappa I\) almost everywhere for one
\(\kappa>0\), then \(\mathcal M\) has closed range and is bounded below.
\end{lemma}

\begin{proof}
The adjoint of the maximal decomposable operator has fibers \(M(x)^*\).
Therefore
\[
 (\overline{\operatorname{Ran}\mathcal M})^\perp
 =\ker\mathcal M^*
 =\int_X^\oplus\ker M(x)^*\dd\mu(x),
\]
which proves (10.31).  The lower-bound assertion is the closed-range
theorem applied fiberwise and then integrated.
\end{proof}

For a countable set \(Q\) of prime powers, apply the lemma to the charged
Gamma operator on \(\bigoplus_{q\in Q}\mathscr K_C\),
\[
 K_{\Gamma,Q}
 =\kappa_\infty I+\bigoplus_{q\in Q}m_\Gamma(A-\log q).
 \tag{10.32}
\]
It is uniformly bounded below by \(\kappa_\infty I\).  If the scalar term
is omitted, the summand \(q\) is multiplication by \(m_\Gamma(\gamma-\log q)\),
which vanishes only at the null set \(\{\log q\}\); hence its adjoint has
zero kernel and, by (10.31), the range is still dense.  Unitary charge transports cannot change kernels
of adjoints or closure of ranges.

\begin{theorem}[Absolutely continuous completion no-go]
\label{thm:ac-no-go}
The Hilbert completion \(\mathscr P_\infty^{\rm cov}\) of the covariant
system (10.20)--(10.29), with its unitary group \(U^{\rm cov}_t\), has
purely absolutely continuous spectrum and no point spectrum.  Consequently,
under Hypothesis (R) for critical-line zeros, every linear map
intertwining \(U_t=e^{-t/2}\vartheta_t\) with \(U^{\rm cov}_t\),
\[
 D:H^1_{\rm CCM}\longrightarrow\mathscr P_\infty^{\rm cov}
 \tag{10.33}
\]
annihilates every resonant eigenclass.  In particular it cannot be a
faithful descent and cannot supply the polarization of
Definition~\ref{def:global-polarization}.
\end{theorem}

\begin{proof}
The measure in (10.20) is equivalent to Lebesgue measure.  Each source
summand of \(\mathscr S_X\) is a copy of \(\mathscr K_C\) with generator
\(A-\ell\), unitarily equivalent to multiplication by \(\gamma-\ell\) on
\(L^2(\R,w_C\,\dd\gamma)\).  For \(v\in\mathscr S_X\),
\[
\begin{aligned}
 g^{\rm cov}(U^{\rm cov}_tv,v)
 &=g_X(U^X_tv,v)\\
 &\quad+\langle e^{itA}K_\Gamma^{-1/2}R_X^{\rm cov}J_Xv,K_\Gamma^{-1/2}R_X^{\rm cov}J_Xv\rangle\\
 &\quad+\langle e^{itA}K_\Gamma^{-1/2}R_X^{\rm cov}v,K_\Gamma^{-1/2}R_X^{\rm cov}v\rangle,
\end{aligned}
\]
using the intertwining relation after (10.24) and \([K_\Gamma,A]=0\).  Each
term is the Fourier transform of a measure absolutely continuous with
respect to Lebesgue measure, so the spectral measure of \(v\) is absolutely
continuous.  Thus \(\bigcup_X\mathscr S_X\) lies in the absolutely continuous
subspace, which is closed; since the union is dense, the whole space is
absolutely continuous, and there are no eigenvectors.

Under Hypothesis (R) for critical-line zeros, for every such zero
\(\rho=1/2+i\gamma\) there is \(u_\rho\ne0\) with
\(U_tu_\rho=e^{it\gamma}u_\rho\).  Equivariance of (10.33)
would make \(Du_\rho\) an eigenvector with the same character.  Since the
target has no point spectrum, \(Du_\rho=0\).  Infinitely many such zeros
exist unconditionally by Hardy's theorem~\cite{Hardy1914}.
\end{proof}

\begin{remark}[Hilbert-target scope]
Theorem~\ref{thm:ac-no-go} does not exclude an equivariant distributional
map into the dual member \(\Phi'\) of a rigged triple
\(\Phi\subset\mathscr P_\infty^{\rm cov}\subset\Phi'\).  Generalized
eigenvectors of an absolutely continuous operator naturally live in such a
dual.  The theorem excludes faithful descent into the Hilbert member, which
is exactly the target of the positive metric (10.27).  A resonant nuclear
construction must therefore retain evaluation vectors and jets in
\(\Phi'\) before asking which combinations admit a positive Hilbert
completion.
\end{remark}

The theorem concerns the single completion \(\mathscr P_\infty^{\rm cov}\).
It does not assert that every conceivable Hilbert completion is absolutely
continuous.  The same argument applies to any completion whose vectors
have absolutely continuous spectral measures on a dense invariant
subspace.

\subsection{Six terminal formulations and their common content}
\label{subsec:six-terminal}

The same sign can be stated in six forms.  For the first two and the last,
the implication to Weil positivity or RH is proved in this paper: the first
by Theorem~\ref{thm:mosco-transfer}, the second by
Theorem~\ref{thm:cofinal-inertia-closure} under the ground-state identity
(on even tests),
the last under Hypothesis (R) by Proposition~\ref{prop:polarization-rh}; the third is
the GNS restatement of positivity of the Weil form; the fourth and fifth
are stated as formulations only and are not developed here:
\begin{center}
\renewcommand{\arraystretch}{1.2}
\begin{tabular}{>{\raggedright\arraybackslash}p{0.27\textwidth}
                >{\raggedright\arraybackslash}p{0.63\textwidth}}
\toprule
Realization & Terminal statement\\
\midrule
Localized signed forms & Lower semicontinuity of the indefinite Weil form
under the global exhaustion, not merely Mosco convergence of its positive
shift.\\
Finite inertia/Schur analysis & The physical prime--Gamma gain exceeds the
anti-shorted deficit, \(\mathfrak g_{M,J,X}>\delta^{\flat}_{M,J,X}\), on a
form-core exhaustion.\\
Fixed Rosati realization & A faithful factorization
\(\tau(f*g^\sharp)=\langle D[f],D[g]\rangle\); by GNS, its existence is
equivalent to positivity of the Rosati/Weil form.\\
Alternative polarization & A source-defined positive Hilbert majorant for
which normalized scaling is unitary and \(\Omega\) is bounded and weakly
nondegenerate.\\
Charged boundary pushout & Closure-faithful equivariant descent from the
complete Tate--Gamma--polar source to the separated CCM cokernel.\\
Geometric degree one & A positive global polarization
\((H^1_{\rm CCM},\Omega,J,g)\) as in Definition~\ref{def:global-polarization}.\\
\bottomrule
\end{tabular}
\end{center}

The first two rows localize the sign in finite matrices.  The third states
the sign as a factorization.  The fourth and fifth try to construct a new
positive metric before comparison with Rosati.  The last states the desired
geometric structure directly.  In each row the unproved
assertion is the terminal statement itself.

This repeated endpoint admits a general formulation.

\begin{proposition}[Deformation-blindness obstruction]
\label{prop:deformation-blindness}
Let \(\mathcal X\) be a class of divisor-bearing objects, let \(\sim\) be a
deformation relation on \(\mathcal X\), and let
\(\kappa:\mathcal X\to\mathbb Z\) be an index depending on the position of
the divisor.  Suppose that a functor \(F:\mathcal X\to\mathcal Y\) is
constant on \(\sim\)-classes.  If \(\kappa\) is not constant on one such
class, then no map \(\phi:F(\mathcal X)\to\mathbb Z\) satisfies
\(\kappa=\phi\circ F\).
\end{proposition}

\begin{proof}
Choose \(x\sim y\) with \(\kappa(x)\ne\kappa(y)\).  Then
\(F(x)=F(y)\), whereas a factorization would give
\(\kappa(x)=\phi(F(x))=\phi(F(y))=\kappa(y)\), a contradiction.
\end{proof}

In the present program, Proposition~\ref{prop:deformation-blindness}
organizes four otherwise separate obstructions.  The table states the
forgetful functor actually used, the relation through which it factors,
and an explicit witness showing that the required index does not factor
through it.

\begin{center}
\small
\renewcommand{\arraystretch}{1.25}
\begin{tabular}{>{\raggedright\arraybackslash}p{0.20\textwidth}
                >{\raggedright\arraybackslash}p{0.26\textwidth}
                >{\raggedright\arraybackslash}p{0.44\textwidth}}
\toprule
Functor \(F\) & Relation \(R\) forgotten by \(F\) & Witness with different
index information\\
\midrule
Half-plane sign class of the Maass--Selberg Cauchy pairing
& Motion of resonance parameters inside \(0<\Re\rho<1\)
& Writing \(w_\rho=(1-\rho)/2\), the kernel is always
\(-1/(w_\rho+\overline{w_\sigma})\), hence negative definite.  Five
on-line resonances give signature \((0,5)\); adjoining an off-line quartet
gives \((0,9)\), still definite, although the Weil off-line index changes.\\
Scalar determinant section
& Equality of determinants after forgetting the sourced Hermitian kernel
and its Krein--Langer divisor
& The Hermitian matrices \(I_2\) and \(-I_2\) both have determinant one,
but negative indices zero and two.  The same ambiguity is why a
Nevanlinna determinant must be augmented by its sourced second variation
and divisor grading.\\
Cyclic Chern character
& Bounded operator homotopy, including scalar shifts of a finite
semilocal Weil operator
& The path from an unshifted finite Weil matrix to its lower-edge-shifted
nonnegative matrix preserves the character, while the lower edge and the
negative eigenvalue count can change (stated, not proved here).\\
Curvature of a positive determinant-line metric
& Motion of a section divisor with the Hermitian line and its metric fixed
& On a disk, \(s_a(z)=z-a\) and \(s_b(z)=z-b\) have different divisors,
but the same curvature \(\partial\bar\partial\phi\) for
\(\|1\|^2=e^{-\phi}\).  Poincar\'e--Lelong records the moved divisor but
metric positivity does not localize it.\\
\bottomrule
\end{tabular}
\end{center}

Theorem~\ref{thm:ac-no-go} is deliberately not a fifth row of this table.
Its proof does not require a deformation witness.  It is a spectral-type
obstruction: under Hypothesis (R) the source has point-spectrum eigenclasses, the Hilbert target
has purely absolutely continuous spectrum, and equivariance forces every
such class to map to zero.

The conclusion is not that every analytic limit is signature-blind.
Spectral flow, for example, can retain an index when a gap and a continuous
Fredholm path are supplied.  The conclusion is narrower and operational:
for the witnesses listed in the table (the cyclic Chern character row is
stated, not proved here), none of the functors can recover the off-line
index because it factors through an equivalence relation under which that index
can change; and the absolutely continuous Hilbert completion fails for the
independent point-versus-continuous spectral reason.

\subsection{Two open programs}
\label{subsec:open-programs}

The preceding construction and no-go leave two mathematically distinct
directions.  Both must be divisor-sensitive from the start; neither may
define its positive metric from the zeros or from the Weil form whose sign
is sought.

\subsubsection{A geometric intersection product and Hodge-index theorem}

The first direction is to construct the missing geometric source of
positivity.  The finite identity (10.7) and the archimedean identity
(10.12) specify its local intersection numbers and determinant fibers.
What is still absent is a global square carrying correspondences, a
relative intersection form, and a primitive Hodge-index theorem.  A
successful object must provide:
\begin{enumerate}
\item classes \(\Gamma_{p^k}\) whose local torsion and incidence recover
\(\Lambda(p^k)/\sqrt{p^k}\);
\item an archimedean fiber whose determinant is (10.12), including the
polar \(H^0/H^2\) page;
\item a Lefschetz or supertrace formula identifying the resulting
intersection distribution with Weil's explicit formula;
\item a primitive intersection inequality whose induced Rosati form on
degree one is positive;
\item compatibility with the scaling correspondences, so that normalized
scaling is unitary on the polarized degree one.
\end{enumerate}
This is the analogue of the role played by correspondences, the Jacobian
and the Hodge index theorem in Weil's proof over finite fields.  The work of
Connes and Consani on the arithmetic curve, its Jacobian and absolute
geometry supplies the natural ambient program
\cite{CCArch,CCJacobian,CCAbsolute}.  Theorems
\ref{thm:root-torsion} and \ref{thm:gamma-polar-det} give concrete local
constraints that any such global theory must satisfy.  They do not supply
the global intersection inequality.

\subsubsection{A divisor-sensitive resonant completion}

The second direction remains analytic or derived.  It must replace the
ordinary direct-integral completion by a topology in which evaluation and
multiplicity jets at a discrete divisor survive before Hilbert reduction.
At minimum one needs a source-defined space \(\mathscr E\), a closed
restriction complex and a map
\[
 H^1_{\rm CCM}\longrightarrow H^1(\mathscr E)
 \tag{10.34}
\]
with the following properties:
\begin{enumerate}
\item (10.34) is faithful on every resonant jet and is equivariant for
scaling;
\item the Tate, Gamma and polar boundary maps act continuously and retain
the exact finite part (10.19);
\item the alternating residue form extends boundedly and nondegenerately;
\item a positive compatible complex structure is constructed from source
data, not from the zero divisor or the positive part of the Weil form;
\item the resulting normalized scale action is unitary.
\end{enumerate}
Possible ambient categories include nuclear analytic spaces, derived
cyclic modules and reproducing-kernel completions.  The category alone is
not a solution: the classical de Branges program shows that bounded point
evaluation can simply re-express the missing Hermite--Biehler sign, and
specific proposed positivity conditions for zeta and \(L\)-functions are
known to fail~\cite{deBranges1968,ConreyLi2000}.  The acceptance test is
therefore exact: the construction must fail for a Davenport--Heilbronn
type function without an Euler product~\cite{DavenportHeilbronn1936}, must preserve the CCM divisor, and
must obtain positivity from arithmetic source structure rather than assume
it.

The two programs are related but not identical.  A geometric intersection
theory could produce the resonant completion as its analytic realization.
Conversely, a successful divisor-sensitive completion might reveal the
one-dimensional shadow of the missing intersection theory.  At present
both positivity theorems remain open.

\section{Conclusion}

The Connes--CCM route to RH, as analysed here, has the following conditional
shape:
\[
\boxed{
\begin{gathered}
\text{finite semilocal Weil self-adjointness}\\
+\ \text{finite critical-line prolate zeta-cycles}\\
+\ \text{Suzuki localized operator realization}\\
+\ \text{a faithful positive global polarization}\\
\Longrightarrow \QW\ge0
\Longrightarrow \RH .
\end{gathered}}
\tag{11.1}
\]
The first three blocks are quoted from the
CCM/Connes/Suzuki programs only in the restricted sense stated in
\S\S2--7: self-adjointness and real finite zeros under the hypothesis of
Theorem~\ref{thm:ccm-input}, the localized realization of
Theorem~\ref{thm:suzuki-input}, while the prolate limit (2.4) and the
spectral identification of the limit are used only as explicit
assumptions.  The finite analysis of \S\S8--9 proves that approximation,
tail control, radical shorting, prime sampling, adaptive gain recovery and
form-core exhaustion reduce the sign on each heat or hybrid row to one
finite inequality, which is not proved; the Section~8 bounds hold under
the standing assumptions (P1)--(P4), and the passage to the Weil form
uses the ground-state identity as a hypothesis.  In that realization
the terminal assertion is the physical surplus
\(\mathfrak g_{M,J,X}>\delta^{\flat}_{M,J,X}\) on the heat or hybrid rows.

Section~\ref{sec:polarization-geometry} then reconstructs the source geometry
behind the same assertion.  At every finite prime, a positive cyclic-root
Laplacian produces the relative torsion \(\log p\) and the incidence
\(p^{-k/2}\), hence the exact coefficient
\(\Lambda(p^k)/\sqrt{p^k}\).  At infinity, the positive Gamma spin and the
polar plane reproduce the completed archimedean determinant.  The common
boundary has an exact matched finite part, and the charge translations
assemble all prime powers equivariantly without removing cross phases.
Theorems~\ref{thm:root-torsion}, \ref{thm:gamma-polar-det} and
\ref{thm:matched-finite-part} are unconditional.

They do not complete (11.1).  Theorem~\ref{thm:ac-no-go} proves that the
ordinary absolutely continuous Hilbert completion of this source object
has no point spectrum and therefore, under Hypothesis (R) for
critical-line zeros, cannot receive the resonant CCM degree one faithfully.  Thus the finite and local polarizations are valid, while
the attempted global descent is not.  This separates a successful local
construction from the still missing global Hodge/Rosati positivity.

The six terminal formulations in \S\ref{subsec:six-terminal} are not six independent missing
lemmas.  They are six formulations of the same missing sign.
Proposition~\ref{prop:deformation-blindness} identifies their recurring
deformation failure: the half-plane sign class, scalar determinants, homotopy classes and curvature
factor through relations which forget the position or signature of the
resonant divisor.  The absolutely continuous Hilbert pushout fails for the
independent spectral reason of Theorem~\ref{thm:ac-no-go}.  Together these
results show, for the witnesses listed in the table (the cyclic Chern
character row being stated, not proved here), that the functors there and the absolutely
continuous completion of Theorem~\ref{thm:ac-no-go} cannot recover the
off-line index.

Two routes remain open.  The geometric route asks for an arithmetic
intersection product and a primitive Hodge-index theorem whose local
fibers are constrained by Theorems~\ref{thm:root-torsion} and
\ref{thm:gamma-polar-det}.  The divisor-sensitive route asks for a
resonant nuclear or derived completion which retains every evaluation jet
before Hilbert reduction and constructs a positive scale-compatible metric
from arithmetic source data.  Either successful construction would supply
the fourth block of (11.1).  Neither construction, nor RH, is proved here.

\begin{thebibliography}{99}

\bibitem{Polya1926}
G.~P\'olya, Bemerkung \"uber die Integraldarstellung der Riemannschen
\(\xi\)-Funktion, \emph{Acta Math.} \textbf{48} (1926), 305--317.

\bibitem{Fuchs1964}
W.~H.~J. Fuchs, On the eigenvalues of an integral equation arising in the
theory of band-limited signals, \emph{J. Math. Anal. Appl.} \textbf{9}
(1964), 317--330.

\bibitem{Slepian1965}
D.~Slepian, Some asymptotic expansions for prolate spheroidal wave
functions, \emph{J. Math. and Phys.} \textbf{44} (1965), 99--140.

\bibitem{Weil1952}
A.~Weil,
\emph{Sur les ``formules explicites'' de la th\'eorie des nombres premiers},
Comm. Sem. Math. Univ. Lund (1952), 252--265.

\bibitem{Connes1999}
A.~Connes,
\emph{Trace formula in noncommutative geometry and the zeros of the
Riemann zeta function},
Selecta Math. (N.S.) \textbf{5} (1999), 29--106.

\bibitem{CCArch}
A.~Connes and C.~Consani,
\emph{Weil positivity and trace formula, the archimedean place},
arXiv:2006.13771 (2020).

\bibitem{CvS2025}
A.~Connes and W.~D. van Suijlekom,
\emph{Quadratic forms, real zeros and echoes of the spectral action},
Comm. Math. Phys. \textbf{406} (2025), no.~12, Paper No.~312.

\bibitem{Connes2026}
A.~Connes,
\emph{The Riemann Hypothesis: Past, Present and a Letter Through Time},
arXiv:2602.04022 (2026).

\bibitem{CCM2025}
A.~Connes, C.~Consani, and H.~Moscovici,
\emph{Zeta spectral triples},
arXiv:2511.22755 (2025).

\bibitem{CCJacobian}
A.~Connes and C.~Consani,
\emph{On the Jacobian of \(\overline{\operatorname{Spec}\mathbb Z}\)},
arXiv:2602.15941 (2026).

\bibitem{CCAbsolute}
A.~Connes and C.~Consani,
\emph{On the Absolute Geometry of \(\operatorname{Spec}\mathbb Z\)},
arXiv:2606.06604 (2026).

\bibitem{Suzuki2025}
M.~Suzuki,
\emph{On the Hilbert space derived from the Weil distribution},
Canadian Journal of Mathematics, FirstView (2025).

\bibitem{SuzukiDirichlet}
M.~Suzuki,
\emph{Detecting zeros of Dirichlet \(L\)-functions via the Riemann
zeta-function},
arXiv:2508.17701 (2025).

\bibitem{SuzukiScrew}
M.~Suzuki,
\emph{Weil's quadratic form via the screw function},
arXiv:2606.09096 (2026).

\bibitem{Burnol2004}
J.-F.~Burnol,
\emph{Two complete and minimal systems associated with the zeros of the
Riemann zeta function},
J. Th\'eor. Nombres Bordeaux \textbf{16} (2004), no.~1, 65--94.

\bibitem{Burnol2007}
J.-F.~Burnol,
\emph{Entrelacement de co-Poisson},
Ann. Inst. Fourier (Grenoble) \textbf{57} (2007), no.~2, 525--602.

\bibitem{Deninger1998}
C.~Deninger,
\emph{Some analogies between number theory and dynamical systems on
foliated spaces},
Doc. Math., Extra Volume ICM 1998, vol.~I, 163--186.

\bibitem{deBranges1968}
L.~de Branges,
\emph{Hilbert Spaces of Entire Functions},
Prentice--Hall, Englewood Cliffs, NJ, 1968.

\bibitem{ConreyLi2000}
J.~B.~Conrey and X.-J.~Li,
\emph{A note on some positivity conditions related to zeta and
\(L\)-functions},
Internat. Math. Res. Notices (2000), no.~18, 929--940.

\bibitem{Hardy1914}
G.~H.~Hardy,
\emph{Sur les z\'eros de la fonction \(\zeta(s)\) de Riemann},
C. R. Acad. Sci. Paris \textbf{158} (1914), 1012--1014.

\bibitem{DavenportHeilbronn1936}
H.~Davenport and H.~Heilbronn,
\emph{On the zeros of certain Dirichlet series},
J. London Math. Soc. \textbf{11} (1936), 181--185, 307--312.

\bibitem{DunsterPSWF}
T.~M. Dunster,
\emph{Asymptotics of prolate spheroidal wave functions},
arXiv:1601.00699v3 (2017).

\bibitem{Krein1947}
M.~G. Krein,
\emph{The theory of self-adjoint extensions of semi-bounded Hermitian
transformations and its applications. I},
Mat. Sb. (N.S.) \textbf{20(62)} (1947), 431--495.

\bibitem{Anderson1971}
W.~N. Anderson, Jr.,
\emph{Shorted operators},
SIAM J. Appl. Math. \textbf{20} (1971), no.~4, 520--525.

\bibitem{AndersonTrapp1975}
W.~N. Anderson, Jr. and G.~E. Trapp,
\emph{Shorted operators. II},
SIAM J. Appl. Math. \textbf{28} (1975), no.~1, 60--71.

\bibitem{Haynsworth1968}
E.~V. Haynsworth,
\emph{Determination of the inertia of a partitioned Hermitian matrix},
Linear Algebra Appl. \textbf{1} (1968), 73--81.

\bibitem{Gantmacher1959}
F.~R. Gantmacher,
\emph{The Theory of Matrices},
Chelsea, New York, 1959.

\bibitem{Andreief1886}
C.~Andr\'eief,
\emph{Note sur une relation entre les int\'egrales d\'efinies des produits
des fonctions},
M\'em. Soc. Sci. Phys. Nat. Bordeaux (3) \textbf{2} (1886), 1--14.

\bibitem{Bhatia1997}
R.~Bhatia,
\emph{Matrix Analysis},
Graduate Texts in Mathematics, vol.~169, Springer, 1997.

\bibitem{Mosco1969}
U.~Mosco,
\emph{Convergence of convex sets and of solutions of variational
inequalities},
Advances in Math. \textbf{3} (1969), 510--585.

\bibitem{KuwaeShioya2003}
K.~Kuwae and T.~Shioya,
\emph{Convergence of spectral structures: a functional analytic theory
and its applications to spectral geometry},
Comm. Anal. Geom. \textbf{11} (2003), no.~4, 599--673.

\bibitem{Post2012}
O.~Post,
\emph{Spectral Analysis on Graph-like Spaces},
Lecture Notes in Mathematics, vol.~2039, Springer, 2012.

\bibitem{SuzukiUemura2014}
K.~Suzuki and T.~Uemura,
\emph{On instability of global path properties of symmetric Dirichlet
forms under Mosco-convergence},
arXiv:1412.0725 (2014).

\end{thebibliography}

\end{document}
